\chapter{Generating Functions as a Computing Tool}

\begin{center}\itshape
A generating function is a clothesline on which we hang up a sequence of numbers for display.\\
\upshape --- Herbert Wilf, \emph{generatingfunctionology}
\end{center}

In 1740 Philippe Naud\'e, a mathematician in Berlin, wrote to Euler with a puzzle: in how many ways can $50$ be written as a sum of seven
different positive integers? Euler did not count them. He wrote down a product of infinitely many factors, $(1+xz)(1+x^2z)(1+x^3z)\cdots$, and
read the answer, $522$, off the coefficient of $x^{50}z^7$. It was the birth of generating functions as a tool for counting.

Volume~3, Chapter~3, used generating functions of probability distributions: $\mathbb E[s^X]$ packs a whole law
into one function, and sums of independent variables become products. This chapter uses the same device as a pure
computing tool. A sequence $(a_n)$ becomes the series $\sum a_nx^n$, and operations on sequences that are awkward
(convolutions, recurrences, sums over all partitions of a set) become operations on functions that are easy
(products, quotients, exponentials). One can then read off a closed form, an exact recurrence, or, as Chapter~13
will show, the asymptotic size of $a_n$.

We first set up \emph{formal} power series, which make these manipulations rigorous without any question of
convergence. Then come three classical uses: linear recurrences, which are exactly the rational generating
functions; \emph{exponential} generating functions, which count labelled structures and obey the exponential
formula; and partitions of integers, with Euler's pentagonal number theorem, whose proof by Franklin is one of
the most charming bijections in mathematics.

\begin{think}
Expand $\frac1{1-x-x^2}=1+x+2x^2+3x^3+5x^4+\cdots$. Where do the Fibonacci numbers come from, and what does the
factorisation of $1-x-x^2$ tell you about them?
\end{think}

\section{Formal power series}

Let $K$ be $\Q$, $\R$ or $\C$. A \emph{formal power series} over $K$ is a sequence $(a_n)_{n\ge0}$ of elements of $K$,
written $A(x)=\sum_{n\ge0}a_nx^n$; nothing is assumed about convergence, and $x$ is not a number. We write
$[x^n]A=a_n$. Sum and product are defined by
\[
\begin{gathered}\sum a_nx^n+\sum b_nx^n=\sum(a_n+b_n)x^n,\\\Bigl(\sum a_nx^n\Bigr)\Bigl(\sum b_nx^n\Bigr)=\sum_n\Bigl(\sum_{k=0}^na_kb_{n-k}\Bigr)x^n.\end{gathered}
\]
With these operations the formal power series form a commutative ring $K[[x]]$ containing the polynomials $K[x]$:
associativity and distributivity are identities between finite sums of coefficients, checked as for polynomials.
The \emph{order} of a non-zero series is the least $n$ with $a_n\ne0$. Since $[x^n](AB)$ involves only $a_0,\dots,a_n$ and
$b_0,\dots,b_n$, the order of $AB$ is the sum of the orders (the lowest coefficients multiply to a non-zero number), so
$K[[x]]$ has no zero divisors.

\begin{proposition}\label{prop:inverse}
A series $A$ is invertible in $K[[x]]$ if and only if $a_0\ne0$. For instance $\frac1{1-cx}=\sum_nc^nx^n$, and more
generally, for $m\ge1$,
\[
\frac1{(1-cx)^m}=\sum_{n\ge0}\binom{n+m-1}{m-1}c^nx^n.
\]
\end{proposition}
\begin{proof}
If $AB=1$ then $a_0b_0=1$. Conversely, if $a_0\ne0$, the equations $\sum_{k=0}^na_kb_{n-k}=[n=0]$ are solved by
$b_0=\frac1{a_0}$ and $b_n=-\frac1{a_0}\sum_{k=1}^na_kb_{n-k}$, recursively. Then $(1-cx)\sum c^nx^n=1$ because the sum
telescopes. The coefficient of $x^n$ in $\bigl(\sum_jc^jx^j\bigr)^m$ is $c^n$ times the number of $m$-tuples of
non-negative integers with sum $n$, which is $\binom{n+m-1}{m-1}$ by the stars-and-bars count (Volume~0, Chapter~8).
\end{proof}

Infinite sums and products are allowed when they make sense coefficientwise: if $A_1,A_2,\dots$ are series whose
orders tend to infinity, each coefficient of $\sum_kA_k$ involves only finitely many $A_k$, so the sum is defined, and
likewise $\prod_k(1+A_k)$, where each coefficient involves the finitely many factors with order at most $n$. In
particular, if $G$ has $g_0=0$, the \emph{composition} $F(G)=\sum_kf_kG^k$ is defined, since $G^k$ has order at least $k$.

The \emph{derivative} is $DA=\sum_{n\ge1}na_nx^{n-1}$. It is linear and satisfies the product rule $D(AB)=(DA)B+A(DB)$,
since both sides have coefficient $\sum_k(k+1)a_{k+1}b_{n-k}+\sum_ka_k(n-k+1)b_{n-k+1}=(n+1)\sum_ka_kb_{n+1-k}$ at
$x^n$. In characteristic $0$, $DA=0$ forces $A$ to be constant, and the chain rule holds:
\[
D\bigl(F(G)\bigr)=(DF)(G)\cdot DG\qquad(g_0=0).
\]
Indeed, for $F=x^k$ it is the product rule iterated, $D(G^k)=kG^{k-1}DG$; by linearity it holds for polynomials; and the
coefficient of $x^n$ on either side only involves the coefficients $f_0,\dots,f_{n+1}$ of $F$, so it holds for $F$
if it holds for the polynomial $f_0+\dots+f_{n+1}x^{n+1}$.

We define $\exp(x)=e^x=\sum_n\frac{x^n}{n!}$ and $\log(1+x)=\sum_{n\ge1}\frac{(-1)^{n-1}x^n}n$, so that $D\exp=\exp$ and
$D\log(1+x)=\frac1{1+x}$.

\begin{proposition}\label{prop:explog}
Let $F$, $G$ be series with $f_0=g_0=0$. Then $\exp(F+G)=\exp(F)\exp(G)$, $\log(1+(e^F-1))=F$, and
$\exp\bigl(\log(1+F)\bigr)=1+F$.
\end{proposition}
\begin{proof}
By the chain rule $D\exp(F)=\exp(F)DF$. The series
\[
H=\exp(F+G)\exp(-F)\exp(-G)
\]
 has
$DH=H\cdot D(F+G-F-G)=0$ by the product rule, so $H$ is constant, equal to its constant term $1$; and
$\exp(-F)\exp(F)=1$ by the same argument with $G=-F$. This gives the first identity. For the second, put
$U=e^F-1$: then $D\log(1+U)=\frac{DU}{1+U}=\frac{e^FDF}{e^F}=DF$, and both sides vanish at $0$. For the third, let
$L=\log(1+F)$, so that $DL=\frac{DF}{1+F}$; then
\[
D\bigl(e^L(1+F)^{-1}\bigr)=e^L(1+F)^{-1}\bigl(DL-\frac{DF}{1+F}\bigr)=0,
\]

and the constant term is $1$.
\end{proof}

\section{Two tools from analysis}

Generating functions are most useful when they also converge, and then two facts about infinite sums are needed again
and again in this volume. Volume~0 used them only in special cases, so we prove them here.

For a family $(a_i)_{i\in I}$ of non-negative numbers indexed by a countable set $I$, let $\sum_{i\in I}a_i\in[0,+\infty]$ be the
supremum of the finite sums $\sum_{i\in F}a_i$ over finite $F\subseteq I$. A complex family is \emph{summable} if
$\sum_{i\in I}|a_i|<\infty$.

\begin{lemma}[summation by packets]\label{lem:packets}
\begin{enumerate}[label=(\alph*)]
\item If $a_i\ge0$, then for every enumeration $i_1,i_2,\dots$ of $I$, $\sum_{k\ge1}a_{i_k}=\sum_{i\in I}a_i$.
\item If $a_i\ge0$ and $I$ is the disjoint union of countably many sets $I_j$, then $\sum_{i\in I}a_i=\sum_j\sum_{i\in I_j}a_i$.
\item If $(a_i)$ is a summable complex family, the series $\sum_ka_{i_k}$ converges absolutely to the same sum $S$ for every
enumeration, and $S=\sum_j\sum_{i\in I_j}a_i$ for every partition into packets $I_j$, each inner sum being itself absolutely convergent.
In particular, for a double sequence with $\sum_{m,n}|a_{mn}|<\infty$, $\sum_m\sum_na_{mn}=\sum_n\sum_ma_{mn}$.
\end{enumerate}
\end{lemma}
\begin{proof}
(a) Each partial sum $\sum_{k\le K}a_{i_k}$ is a finite sum, hence at most $\sum_Ia_i$; and each finite $F\subseteq I$ is contained in
$\{i_1,\dots,i_K\}$ for some $K$, so $\sum_Fa_i\le\sum_{k\le K}a_{i_k}$. The two suprema are equal.

(b) A finite $F\subseteq I$ meets finitely many $I_j$, and $\sum_Fa_i=\sum_j\sum_{F\cap I_j}a_i\le\sum_j\sum_{I_j}a_i$; so
$\sum_Ia_i\le\sum_j\sum_{I_j}a_i$. Conversely, for finitely many indices $j\le J$ and finite sets $F_j\subseteq I_j$,
$\sum_{j\le J}\sum_{F_j}a_i=\sum_{\bigcup F_j}a_i\le\sum_Ia_i$; taking suprema over the $F_j$, then over $J$, gives the reverse
inequality.

(c) Write $a_i=b_i-c_i+i(d_i-e_i)$ with $b_i=\max(\operatorname{Re}a_i,0)$, $c_i=\max(-\operatorname{Re}a_i,0)$ and similarly $d_i$, $e_i$
for the imaginary part. These four families are non-negative with sums at most $\sum|a_i|<\infty$, so (a) and (b) apply to each,
with finite values, and the statements for $(a_i)$ follow by linearity of convergent series. Absolute convergence of each
enumeration and of each inner sum holds by (a) and (b) applied to $(|a_i|)$.
\end{proof}

A sequence of functions $f_n$ on a set $E$ \emph{converges uniformly} to $f$ if $\sup_E|f_n-f|\to0$.

\begin{lemma}[uniform convergence]\label{lem:uniform}
Let $E$ be an interval.
\begin{enumerate}[label=(\alph*)]
\item A uniform limit of continuous functions on $E$ is continuous.
\item \emph{Weierstrass test}: if $|u_n(x)|\le M_n$ on $E$ with $\sum M_n<\infty$, then $\sum u_n$ converges uniformly on $E$.
\item If $f_n\to f$ uniformly on $[a,b]$ and the $f_n$ are continuous, then $\int_a^bf_n\to\int_a^bf$. So a uniformly convergent
series of continuous functions can be integrated term by term.
\item If the $f_n$ are of class $C^1$ on $E$, $f_n(x_0)$ converges for some $x_0\in E$, and $f_n'\to g$ uniformly on $E$, then $f_n$
converges, uniformly on every bounded subinterval, to a function $f$ of class $C^1$ with $f'=g$. So a series of $C^1$ functions can
be differentiated term by term when the differentiated series converges uniformly.
\end{enumerate}
\end{lemma}
\begin{proof}
(a) Let $x\in E$ and $\varepsilon>0$. Choose $n$ with $\sup|f_n-f|\le\varepsilon$, then $\delta>0$ with $|f_n(y)-f_n(x)|\le\varepsilon$ for
$|y-x|\le\delta$. Then $|f(y)-f(x)|\le|f(y)-f_n(y)|+|f_n(y)-f_n(x)|+|f_n(x)-f(x)|\le3\varepsilon$.
(b) The series converges absolutely at each $x$, and the tail satisfies $\bigl|\sum_{k>n}u_k(x)\bigr|\le\sum_{k>n}M_k\to0$ uniformly in $x$.
(c) $f$ is continuous by (a), hence integrable, and $\bigl|\int_a^b(f_n-f)\bigr|\le(b-a)\sup|f_n-f|\to0$.
(d) By (a), $g$ is continuous. By the fundamental theorem of calculus, $f_n(x)=f_n(x_0)+\int_{x_0}^xf_n'$. By (c) the right side
converges to $f(x)=\ell+\int_{x_0}^xg$, where $\ell=\lim f_n(x_0)$, uniformly for $x$ in a bounded subinterval $J$, since
$\bigl|\int_{x_0}^x(f_n'-g)\bigr|\le|x-x_0|\sup_E|f_n'-g|$. And $f'=g$ by the fundamental theorem again.
\end{proof}

\begin{remark}
When the series converge, formal computations agree with the functions. If $\sum a_nz^n$ and $\sum b_nz^n$ converge
absolutely for $|z|<r$, their Cauchy product converges to the product of the sums (Volume~0, Chapter~5). And a series
with positive radius of convergence is determined by its sum near $0$: if $\sum c_nz^n=0$ for all small $z$ and $c_m$ is
the first non-zero coefficient, dividing by $z^m$ and letting $z\to0$ gives $c_m=0$. So an identity between functions
that are sums of power series near $0$ gives an identity of formal series, and conversely. We use this freely below.
\end{remark}

\section{Rational series and linear recurrences}

\begin{theorem}\label{thm:rational}
Let $a_1,\dots,a_d\in\C$ with $a_d\ne0$, and $Q(x)=1-a_1x-\dots-a_dx^d$. For a sequence $(u_n)$ with generating
function $U=\sum u_nx^n$, the following are equivalent:
\begin{enumerate}[label=(\roman*)]
\item $u_{n+d}=a_1u_{n+d-1}+\dots+a_du_n$ for every $n\ge0$;
\item $U=\frac PQ$ for a polynomial $P$ of degree less than $d$.
\end{enumerate}
In that case, if $Q(x)=\prod_{i=1}^s(1-\alpha_ix)^{m_i}$ with distinct $\alpha_i$, there are polynomials $p_i$ of degree less
than $m_i$ such that $u_n=\sum_{i=1}^sp_i(n)\alpha_i^n$ for all $n\ge0$.
\end{theorem}
\begin{proof}
The coefficient of $x^{n+d}$ in $QU$ is $u_{n+d}-a_1u_{n+d-1}-\dots-a_du_n$. So (i) says exactly that $P=QU$ has no
terms of degree $\ge d$, which is (ii) since $Q$ is invertible.

For the explicit form, note first that $Q(0)=1$ and that $Q$ factors over $\C$ as $\prod(1-\alpha_ix)^{m_i}$: the $\alpha_i$
are the inverses of the roots of $Q$ (equivalently, the roots of $y^d-a_1y^{d-1}-\dots-a_d$, all non-zero since $a_d\ne0$).
We show by induction on the number $s$ of distinct factors that $\frac PQ$, with $\deg P<\deg Q$, is a sum
$\sum_i\frac{P_i}{(1-\alpha_ix)^{m_i}}$ with $\deg P_i<m_i$. For $s=1$ there is nothing to do. For $s\ge2$, write $Q=Q_1Q_2$
with $Q_1=(1-\alpha_1x)^{m_1}$; $Q_1$ and $Q_2$ have no common root, so they are coprime, and since $\C[x]$ is a principal
ideal domain (Volume~2, Chapter~6) there are polynomials $A$, $B$ with $AQ_1+BQ_2=1$. Then
$\frac PQ=\frac{PA}{Q_2}+\frac{PB}{Q_1}$. Divide $PB$ by $Q_1$: $PB=CQ_1+R_1$ with $\deg R_1<m_1$, so
$\frac PQ=\frac{R_1}{Q_1}+\frac{PA+CQ_2}{Q_2}$. Write $R_2=PA+CQ_2$. Multiplying by $Q$ gives $P=R_1Q_2+R_2Q_1$, and since
$\deg P<\deg Q$ and $\deg(R_1Q_2)<\deg Q$, also $\deg(R_2Q_1)<\deg Q$, so $\deg R_2<\deg Q_2$. The induction hypothesis
applies to $\frac{R_2}{Q_2}$.

Finally, a polynomial $P_i$ of degree less than $m$ can be written in powers of $1-\alpha x$ (substitute
$x=\frac{1-(1-\alpha x)}\alpha$ and expand): $P_i=\sum_{j=0}^{m-1}c_j(1-\alpha x)^j$. So
$\frac{P_i}{(1-\alpha x)^m}=\sum_{j=1}^m\frac{c_{m-j}}{(1-\alpha x)^j}$, and by Proposition~\ref{prop:inverse} its coefficient
of $x^n$ is $\alpha^n\sum_jc_{m-j}\binom{n+j-1}{j-1}$, a polynomial in $n$ of degree less than $m$ times $\alpha^n$.
\end{proof}

\begin{example}[Fibonacci]\label{ex:fibonacci}
$F_0=0$, $F_1=1$, $F_{n+2}=F_{n+1}+F_n$. Here $Q=1-x-x^2$ and $P=QU$ has coefficients $F_0=0$ and $F_1-F_0=1$, so
$U=\frac x{1-x-x^2}$. With $\varphi=\frac{1+\sqrt5}2$ and $\psi=\frac{1-\sqrt5}2$, $1-x-x^2=(1-\varphi x)(1-\psi x)$, and
\[
\begin{gathered}\frac x{(1-\varphi x)(1-\psi x)}=\frac1{\sqrt5}\Bigl(\frac1{1-\varphi x}-\frac1{1-\psi x}\Bigr),\\F_n=\frac{\varphi^n-\psi^n}{\sqrt5}.\end{gathered}
\]
This is Binet's formula; Exercise~\ref{ex:binet} draws consequences.
\end{example}

The method also handles recurrences that are not of constant-coefficient type, as long as the generating
function satisfies an equation one can solve. Two examples from the analysis of algorithms show the range.

\begin{example}[merge sort]\label{ex:mergesort}
Merge sort splits a list of $n=2^m$ items in two halves, sorts each, and merges them with at most $n$ comparisons. Its
worst-case cost $T(n)$ satisfies $T(n)=2T(\frac n2)+n$, so $t_m=T(2^m)$ satisfies $t_m=2t_{m-1}+2^m$. Multiplying by
$y^m$ and summing over $m\ge1$, the series $\tau(y)=\sum t_my^m$ satisfies $\tau-t_0=2y\tau+\frac{2y}{1-2y}$, so
\[
\tau(y)=\frac{t_0}{1-2y}+\frac{2y}{(1-2y)^2},\qq t_m=t_02^m+m2^m,
\]
by Proposition~\ref{prop:inverse}. In terms of $n$: $T(n)=n\log_2n+T(1)n$.
\end{example}

\begin{example}[quicksort]\label{ex:quicksortexact}
Randomised quicksort (Volume~3, Exercise~2.3) chooses a uniform pivot among $n$ items, compares it with the $n-1$
others, and recurses on both sides. Its expected number of comparisons satisfies $C_0=0$ and
\[
C_n=n-1+\frac2n\sum_{k=0}^{n-1}C_k\qquad(n\ge1).
\]
Multiply by $nx^n$ and sum: with $C(x)=\sum C_nx^n$, and $\sum_nn(n-1)x^n=\frac{2x^2}{(1-x)^3}$ (differentiate
$\frac1{1-x}$ twice), this becomes $xC'=\frac{2x^2}{(1-x)^3}+\frac{2xC}{1-x}$, a linear differential equation. Dividing
by $x$, the series $(1-x)^2C$ has derivative $(1-x)^2C'-2(1-x)C=\frac{2x}{1-x}$, so, both sides vanishing at $0$,
\[
\begin{gathered}(1-x)^2C=2\bigl(-\log(1-x)-x\bigr),\\C=\frac2{(1-x)^2}\log\frac1{1-x}-\frac{2x}{(1-x)^2}.\end{gathered}
\]
Since $[x^n]\frac1{(1-x)^2}\log\frac1{1-x}=\sum_{k=1}^n\frac{n+1-k}k=(n+1)H_n-n$ and $[x^n]\frac x{(1-x)^2}=n$,
\[
\begin{aligned}C_n&=2(n+1)H_n-4n\eqbr=2n\ln n+(2\gamma-4)n+O(\ln n),\end{aligned}
\]
using $H_n=\ln n+\gamma+O(\frac1n)$ (Volume~0, Chapter~5, where $\gamma\approx0.5772$ is Euler's constant).
\end{example}

\section{Exponential generating functions}

When the objects counted are built on a set of labels $\{1,\dots,n\}$ (permutations, set partitions, graphs), the right
generating function is the \emph{exponential} one, $\hat A(x)=\sum_na_n\frac{x^n}{n!}$. The reason is the binomial
coefficient in its product:
\[
\hat A\hat B=\sum_n\Bigl(\sum_{k=0}^n\binom nka_kb_{n-k}\Bigr)\frac{x^n}{n!}.
\]
If $a_k$ counts the structures of a kind $\mathcal A$ on a $k$-element set and $b_k$ those of a kind $\mathcal B$, then
$\sum_k\binom nka_kb_{n-k}$ counts the ways to split an $n$-set into two parts $S$, $T$ (with $S$ first) and put an
$\mathcal A$-structure on $S$ and a $\mathcal B$-structure on $T$. Only the size of a set matters for the number of
structures on it, which is what ``labelled'' means.

\begin{theorem}[exponential formula]\label{thm:expformula}
Let $c_k$ be the number of ``connected'' structures of some kind on a $k$-element set ($k\ge1$), and let $h_n$ be the
number of ways to partition an $n$-element set into blocks and put a connected structure on each block ($h_0=1$).
Then, with $\hat C=\sum_{k\ge1}c_k\frac{x^k}{k!}$,
\[
\sum_{n\ge0}h_n\frac{x^n}{n!}=\exp\bigl(\hat C(x)\bigr).
\]
\end{theorem}
\begin{proof}
By the product rule for exponential series, iterated, $[\frac{x^n}{n!}]\hat C^k$ is the number of ways to choose an
\emph{ordered} sequence $(S_1,\dots,S_k)$ of non-empty disjoint sets with union $\{1,\dots,n\}$ and a connected
structure on each. The blocks of a set partition are distinct non-empty sets, so each unordered choice with $k$
blocks appears exactly $k!$ times in this count. Hence the number of structures with exactly $k$ blocks is
$[\frac{x^n}{n!}]\frac{\hat C^k}{k!}$, and summing over $k$ gives $[\frac{x^n}{n!}]\exp(\hat C)$.
\end{proof}

\begin{example}\label{ex:expformula}
\begin{enumerate}[label=(\alph*)]
\item \emph{Permutations.} A permutation is a set of cycles, and there are $(k-1)!$ cyclic orders on a $k$-set. So
$\hat C=\sum_k\frac{x^k}k=\log\frac1{1-x}$ and $\exp(\hat C)=\frac1{1-x}$ by Proposition~\ref{prop:explog}: there are $n!$
permutations, as there should be.
\item \emph{Derangements}, permutations without fixed points: cycles of length at least $2$, so
$\hat D=\exp\bigl(\log\frac1{1-x}-x\bigr)=\frac{e^{-x}}{1-x}$, and $D_n=n!\sum_{k=0}^n\frac{(-1)^k}{k!}$, which is the integer
nearest to $\frac{n!}e$ for $n\ge1$.
\item \emph{Set partitions.} Each block carries one structure, so $c_k=1$ and $\hat C=e^x-1$. The Bell numbers $b_n$,
which count all partitions of an $n$-set, have $\sum b_n\frac{x^n}{n!}=e^{e^x-1}$; they are $1,1,2,5,15,52,203,877,\dots$
Marking each block with a variable $u$ gives $\sum_{n,k}\stirlingii nk u^k\frac{x^n}{n!}=e^{u(e^x-1)}$, the generating
function of the Stirling numbers of Chapter~1.
\item \emph{Involutions}, permutations with $\sigma^2=\mathrm{id}$: cycles of length $1$ or $2$, so their exponential
generating function is $e^{x+x^2/2}$ (Exercise~\ref{ex:involutions}).
\end{enumerate}
\end{example}

The formula also runs backwards: knowing all structures, the logarithm counts the connected ones. Exercise
\ref{ex:connected} counts connected graphs this way, a problem with no simple direct solution.

\section{Partitions of integers}

A \emph{partition} of $n\ge0$ is a way of writing $n$ as a sum of positive integers, the order of the terms not
mattering: $4=3+1=2+2=2+1+1=1+1+1+1$, so $p(4)=5$. We write a partition as $\lambda_1\ge\lambda_2\ge\dots\ge\lambda_r$, and draw
its \emph{Ferrers diagram}: $r$ rows of dots, of lengths $\lambda_1,\dots,\lambda_r$, left aligned (Figure~\ref{fig:ferrers}).
Reflecting the diagram in its diagonal gives the \emph{conjugate} partition, whose parts are the column lengths.

\begin{figure}[ht]
\centering
\fig{ch02-ferrers}
\caption{The Ferrers diagram of $6+4+4+2+1$ and of its conjugate $5+4+3+3+1+1$. Conjugation shows that the number of
partitions of $n$ into at most $k$ parts equals the number of partitions of $n$ into parts at most $k$.}
\label{fig:ferrers}
\end{figure}

\begin{proposition}[Euler]\label{prop:eulerproduct}
In $\Z[[x]]$,
\[
\begin{gathered}\sum_{n\ge0}p(n)x^n=\prod_{k\ge1}\frac1{1-x^k},\qq\prod_{k\ge1}(1+x^k)=\prod_{k\ge1}\frac1{1-x^{2k-1}}.\end{gathered}
\]
The second identity says that the partitions of $n$ into distinct parts are as many as the partitions of $n$ into odd
parts.
\end{proposition}
\begin{proof}
The products make sense since $\frac1{1-x^k}-1$ has order $k$. Expanding $\prod_k(1+x^k+x^{2k}+\cdots)$, the coefficient of
$x^n$ counts the ways to choose $m_k\ge0$ for each $k$, almost all zero, with $\sum_km_kk=n$: a partition with $m_k$ parts
equal to $k$. Similarly $\prod(1+x^k)$ counts partitions into distinct parts, and $\prod\frac1{1-x^{2k-1}}$ partitions into odd
parts. For the identity, $1+x^k=\frac{1-x^{2k}}{1-x^k}$, so for each $N$,
\[
\begin{aligned}\prod_{k\le2N}(1+x^k)&=\frac{\prod_{k\le2N}(1-x^{2k})}{\prod_{k\le2N}(1-x^k)}\eqbr=\frac{\prod_{N<k\le2N}(1-x^{2k})}{\prod_{k\le N}(1-x^{2k-1})}.\end{aligned}
\]
Comparing coefficients of $x^n$ for $n<2N$, where the factors $1-x^{2k}$ with $k>N$ and $1-x^m$ with $m\ge2N$ do not
matter, gives the identity.
\end{proof}

Now the surprise. Expand $\prod_k(1-x^k)$, the inverse of the partition series:
\[
\begin{aligned}\prod_{k\ge1}(1-x^k)&=1-x-x^2+x^5+x^7\eqbr\quad-x^{12}-x^{15}+x^{22}+x^{26}-\cdots\end{aligned}
\]
Almost all coefficients vanish, the others are $\pm1$, and the exponents $1,2,5,7,12,15,\allowbreak22,26,\dots$ are the
\emph{pentagonal numbers} $\frac{j(3j-1)}2$ for $j=\pm1,\pm2,\dots$

\begin{theorem}[Euler's pentagonal number theorem]\label{thm:pentagonal}
\[
\begin{aligned}\prod_{k\ge1}(1-x^k)&=\sum_{j\in\Z}(-1)^jx^{j(3j-1)/2}\\&=1+\sum_{r\ge1}(-1)^rx^{r(3r-1)/2}\eqbr\quad+\sum_{r\ge1}(-1)^rx^{r(3r+1)/2}.\end{aligned}
\]
\end{theorem}
\begin{proof}[Proof (Franklin, 1881)]
Expanding the product, each partition of $n$ into distinct parts contributes $(-1)^r$, where $r$ is its number of
parts. So $[x^n]\prod(1-x^k)=e_n-o_n$, where $e_n$ and $o_n$ count the partitions of $n$ into an even, respectively odd,
number of distinct parts. We pair most of them off.

For a partition $\lambda_1>\dots>\lambda_r$ into distinct parts ($r\ge1$), let $s=\lambda_r$ be its smallest part and let its
\emph{slope} $\sigma$ be the largest $i$ such that $\lambda_1,\dots,\lambda_i$ are consecutive integers, $\lambda_j=\lambda_1-j+1$ for
$j\le i$ (Figure~\ref{fig:franklin}). Call $\lambda$ \emph{exceptional} if $\sigma=r$ and $s\in\{\sigma,\sigma+1\}$. Define:
\begin{itemize}
\item move $A$, when $s\le\sigma$: remove the smallest part and add $1$ to each of the $s$ largest parts;
\item move $B$, when $s>\sigma$: subtract $1$ from each of the $\sigma$ largest parts and add a new part equal to $\sigma$.
\end{itemize}
Let $\mathcal A$ be the non-exceptional partitions with $s\le\sigma$ and $\mathcal B$ those with $s>\sigma$.

\emph{$A$ maps $\mathcal A$ to $\mathcal B$.} Let $\lambda\in\mathcal A$. Then $s\le r-1$: either $\sigma<r$ and $s\le\sigma$, or $\sigma=r$
and $s<\sigma$ since $\lambda$ is not exceptional. So after removing $\lambda_r$ there are at least $s$ parts left, and the result
$\mu$ has parts $\lambda_1+1>\dots>\lambda_s+1>\lambda_{s+1}>\dots>\lambda_{r-1}$, which are distinct since $\lambda_s+1>\lambda_{s+1}$. Its $s$
largest parts are consecutive, and the next one (if any) satisfies $(\lambda_s+1)-\lambda_{s+1}\ge2$: if $s<\sigma$ because
$\lambda_{s+1}=\lambda_s-1$, and if $s=\sigma$ because the slope of $\lambda$ stops at $\sigma$. So $\sigma(\mu)=s$, while the smallest part
of $\mu$ is at least $\lambda_{r-1}>s$. Hence $s(\mu)>\sigma(\mu)$. And $\mu$ is not exceptional: that would need $\sigma(\mu)=r-1$,
so $s=r-1$, and $s(\mu)\in\{\sigma(\mu),\sigma(\mu)+1\}$; since $s(\mu)>\sigma(\mu)$, this means $s(\mu)=\lambda_{r-1}+1=\sigma(\mu)+1=r$, so $\lambda_{r-1}=r-1=s=\lambda_r$, which is impossible.

\emph{$B$ maps $\mathcal B$ to $\mathcal A$.} Let $\mu\in\mathcal B$ with $r$ parts, slope $\sigma$ and smallest part $s>\sigma$. The
result $\nu$ has parts $\mu_1-1>\dots>\mu_\sigma-1$, then $\mu_{\sigma+1},\dots,\mu_r$ if $\sigma<r$, then $\sigma$. These are distinct: if
$\sigma<r$, then $\mu_{\sigma+1}\le\mu_\sigma-2$ by maximality of the slope, and $\mu_r=s>\sigma$; if $\sigma=r$, then $\mu_r-1=s-1>\sigma$
because $s\ne\sigma+1$ ($\mu$ is not exceptional). So $\nu$ is a partition into $r+1$ distinct parts, with smallest part $\sigma$
and slope at least $\sigma$, since its $\sigma$ largest parts are consecutive. So $s(\nu)\le\sigma(\nu)$. If $\nu$ were exceptional
we would have $\sigma(\nu)=r+1$ and $s(\nu)\in\{\sigma(\nu),\sigma(\nu)+1\}$, so $\sigma\ge r+1$, which is impossible. So $\nu\in\mathcal A$.

\emph{$A$ and $B$ are inverse.} For $\lambda\in\mathcal A$, the partition $\mu=A(\lambda)$ has slope $s$, so $B$ subtracts $1$ from its
$s$ largest parts and adds a part $s$: this is $\lambda$. For $\mu\in\mathcal B$, the partition $\nu=B(\mu)$ has smallest part $\sigma$,
so $A$ removes it and adds $1$ to the $\sigma$ largest parts: this is $\mu$.

So $A$ is a bijection from $\mathcal A$ onto $\mathcal B$ that changes the number of parts by one, and the non-exceptional
partitions cancel in $e_n-o_n$. An exceptional partition with $r$ parts is $(2r-1)+(2r-2)+\dots+r$, of sum $\frac{r(3r-1)}2$
(when $s=\sigma=r$), or $2r+(2r-1)+\dots+(r+1)$, of sum $\frac{r(3r+1)}2$ (when $s=r+1$). These numbers are all distinct for
$r\ge1$, so each $n\ge1$ has at most one exceptional partition, contributing $(-1)^r$, and $n=0$ has the empty partition,
contributing $1$. Writing $j=-r$ in the second family, $\frac{r(3r+1)}2=\frac{j(3j-1)}2$, gives the sum over $j\in\Z$.
\end{proof}

\begin{figure}[ht]
\centering
\fig{ch02-franklin}
\caption{Franklin's involution. Left: $7+6+5+2$, with slope $\sigma=3$ (coral) and smallest part $s=2$ (charcoal). Since
$s\le\sigma$, move $A$ puts the smallest part onto the slope: $8+7+5$, whose slope is $2$ and smallest part $5$. Move $B$
takes the slope back off. The two partitions have $4$ and $3$ parts, so they cancel in $\prod(1-x^k)$.}
\label{fig:franklin}
\end{figure}

\begin{corollary}[Euler's recurrence]\label{cor:partitionrec}
For $n\ge1$, with $p(m)=0$ for $m<0$,
\[
\begin{aligned}p(n)&=p(n-1)+p(n-2)\eqbr\quad-p(n-5)-p(n-7)\eqbr\quad+p(n-12)+p(n-15)-\cdots\end{aligned}
\]
the signs following the pattern $++--$ and the arguments $n-\frac{j(3j-1)}2$ for $j=1,-1,2,-2,\dots$
\end{corollary}
\begin{proof}
Multiply $\sum p(n)x^n$ by $\prod(1-x^k)$: the product is $1$, so for $n\ge1$,
$\sum_{j\in\Z}(-1)^jp\bigl(n-\frac{j(3j-1)}2\bigr)=0$. The term $j=0$ is $p(n)$; move the others to the right.
\end{proof}

This is how MacMahon computed $p(n)$ up to $n=200$ by hand in 1918, finding $p(200)=3\,972\,999\,029\,388$; the
numbers grow like $e^{\pi\sqrt{2n/3}}$, as Chapter~4 proves, and the recurrence needs only about $\sqrt n$ terms
for each new value.

\section{Exercises}

\begin{exercise}[$\star$]\label{ex:binet}
Deduce from Binet's formula that $F_n$ is the integer nearest to $\frac{\varphi^n}{\sqrt5}$, that
$\frac{F_{n+1}}{F_n}\to\varphi$, and that $F_{n-1}F_{n+1}-F_n^2=(-1)^n$.
\end{exercise}

\begin{exercise}[$\star$]\label{ex:change}
In how many ways can one pay $1$ euro with coins of $1,2,5,10,20,50$ cents and $1$ euro? Explain why multiplying a
series by $\frac1{1-x^c}$ is a running sum with step $c$, and turn this into an algorithm.
\end{exercise}

\begin{exercise}[$\star\star$]\label{ex:dobinski}
Show that the Bell numbers satisfy $b_n=\frac1e\sum_{k\ge0}\frac{k^n}{k!}$ (Dobi\'nski's formula), and that
$b_{n+1}=\sum_{k=0}^n\binom nkb_k$.
\end{exercise}

\begin{exercise}[$\star\star$]\label{ex:involutions}
Let $t_n$ be the number of involutions of an $n$-set. Show that $t_{n+1}=t_n+nt_{n-1}$ in two ways: from the generating
function $e^{x+x^2/2}$, and directly.
\end{exercise}

\begin{exercise}[$\star\star$]\label{ex:connected}
There are $2^{\binom n2}$ graphs on the vertex set $\{1,\dots,n\}$. Show that the numbers $c_n$ of \emph{connected}
ones satisfy $\sum_{n\ge1}c_n\frac{x^n}{n!}=\log\sum_{n\ge0}2^{\binom n2}\frac{x^n}{n!}$, and compute $c_1,\dots,c_5$.
\end{exercise}

\begin{exercise}[$\star$]\label{ex:motzkin}
A \emph{Motzkin path} of length $n$ goes from $(0,0)$ to $(n,0)$ with steps $(1,1)$, $(1,-1)$, $(1,0)$, never below the
axis. Show that their generating function satisfies $M=1+xM+x^2M^2$, deduce $M=\frac{1-x-\sqrt{1-2x-3x^2}}{2x^2}$, and
compute the first values $1,1,2,4,9,21,51,\dots$
\end{exercise}

\begin{exercise}[$\star\star\star$]\label{ex:sigma}
Let $\sigma(n)$ be the sum of the divisors of $n$. Taking the logarithmic derivative of the pentagonal theorem, show
that for $n\ge1$
\[
\begin{aligned}\sigma(n)&=\sigma(n-1)+\sigma(n-2)\eqbr\quad-\sigma(n-5)-\sigma(n-7)\eqbr\quad+\sigma(n-12)+\sigma(n-15)-\cdots,\end{aligned}
\]
with the same signs as in Corollary~\ref{cor:partitionrec}, the convention $\sigma(0)=n$ for the term $\sigma(n-n)$
if $n$ is pentagonal, and $\sigma(m)=0$ for $m<0$. Euler found this in 1747 and called it ``a most extraordinary law of the
numbers relating to the sum of their divisors''.
\end{exercise}

\begin{exercise}[$\star\star$]\label{pb:prouhet}
Let $s(n)$ be the number of ones in the binary expansion of $n$, and $t_n=(-1)^{s(n)}$.
\begin{enumerate}[label=(\alph*)]
\item Let $N\ge1$. Show that $\sum_{n=0}^{2^N-1}t_n\,n^k=0$ for $k=0,1,\dots,N-1$.
\item Show that for $k=N$ the sum equals $(-1)^NN!\,2^{N(N-1)/2}$.
\item Split $\{0,1,\dots,15\}$ into two sets of eight numbers with the same sum, the same sum of squares and the same sum
of cubes.
\end{enumerate}
\end{exercise}

\begin{exercise}[$\star\star\star$]\label{pb:phipq}
\begin{enumerate}[label=(\alph*)]
\item Let $p<q$ be primes. Show that every coefficient of $\Phi_{pq}$ is $-1$, $0$ or $1$, and that the non-zero
coefficients alternate between $1$ and $-1$.
\item Show that the coefficient of $x^7$ in $\Phi_{105}$ is $-2$.
\end{enumerate}
\end{exercise}

\begin{exercise}[$\star\star\star$]\label{pb:commuting}
Let $h_k(n)$ be the number of $k$-tuples $(\sigma_1,\dots,\sigma_k)$ of permutations of $\{1,\dots,n\}$ that commute pairwise,
with $h_k(0)=1$. Show that
\[
\sum_{n\ge0}h_2(n)\frac{x^n}{n!}=\prod_{j\ge1}\frac1{1-x^j},\qquad\sum_{n\ge0}h_3(n)\frac{x^n}{n!}=\prod_{j\ge1}\Bigl(\frac1{1-x^j}\Bigr)^{\sigma(j)},
\]
where $\sigma(j)$ is the sum of the divisors of $j$. In particular $h_2(n)=n!\,p(n)$, with $p(n)$ the number of partitions of $n$.
\end{exercise}

\begin{exercise}[$\star\star$]\label{pb:roots}
Let $a_n(k)$ be the number of $\sigma\in\mathfrak S_n$ with $\sigma^k=e$, with $a_0(k)=1$.
\begin{enumerate}[label=(\alph*)]
\item Show that $\sum_na_n(k)\frac{x^n}{n!}=\exp\bigl(\sum_{d\mid k}\frac{x^d}d\bigr)$.
\item Let $p$ be prime. Show that $a_n(p)=a_{n-1}(p)+\frac{(n-1)!}{(n-p)!}a_{n-p}(p)$ for $n\ge p$, and deduce that $p$ divides
$a_n(p)$ for every $n\ge p$.
\item Show that $a_n(p)\equiv p\pmod{p(p-1)}$ for $n\ge p$.
\end{enumerate}
\end{exercise}

\begin{exercise}[$\star\star\star$]\label{pb:squares}
Let $s_n$ be the number of permutations of $\{1,\dots,n\}$ that are squares, $\sigma=\tau^2$ for some $\tau\in\mathfrak S_n$.
\begin{enumerate}[label=(\alph*)]
\item Show that $\sigma$ is a square if and only if, for every even $\ell$, $\sigma$ has an even number of cycles of length $\ell$.
\item Show that
\[
\sum_{n\ge0}s_n\frac{x^n}{n!}=\sqrt{\frac{1+x}{1-x}}\ \prod_{m\ge1}\cosh\frac{x^{2m}}{2m}.
\]
\item Deduce that $s_{2m+1}=(2m+1)s_{2m}$, and that the probability that a random permutation is a square is at least that of
Problem~\ref{pb:oddcycles}(b).
\end{enumerate}
\end{exercise}

\begin{exercise}[$\star\star$]\label{pb:randpoly}
\emph{Random polynomials, random permutations.} Let $f$ be a uniformly random monic polynomial of degree $n$ over $\mathbb F_q$,
and $M(j,q)$ the number of monic irreducible polynomials of degree $j$.
\begin{enumerate}[label=(\alph*)]
\item Show that, for $n\ge2$, $f$ is squarefree with probability exactly $1-\frac1q$.
\item Show that $f$ is irreducible with probability $\frac1n+\varepsilon$ with $|\varepsilon|\le\frac2n q^{-n/2}$.
\item For a partition $\lambda$ of $n$ with $m_j$ parts equal to $j$, show that the probability that $f$ factors into irreducibles of
degrees given by $\lambda$ (with multiplicity) tends, as $q\to\infty$, to $1/z_\lambda$, where $z_\lambda=\prod_jj^{m_j}m_j!$.
Interpret $1/z_\lambda$.
\end{enumerate}
\end{exercise}

\begin{exercise}[$\star\star$]\label{pb:parts123}
\emph{Parts $1$, $2$, $3$.} Let $a(n)$ be the number of partitions of $n$ into parts equal to $1$, $2$ or $3$.
\begin{enumerate}[label=(\alph*)]
\item Show that $a(n)-a(n-6)=n$ for $n\ge6$.
\item Deduce that $a(n)$ is the integer nearest to $\frac{(n+3)^2}{12}$ for every $n\ge0$.
\item Show that the number of triangles with integer sides and perimeter $2m+3$, up to congruence, is $a(m)$, and that for
perimeter $2m$ it is $a(m-3)$.
\end{enumerate}
\end{exercise}

\begin{exercise}[$\star$]\label{pb:selfconj}
\emph{Parity by reflection.} A partition is \emph{self-conjugate} if it equals its conjugate. Let $sc(n)$ be the number of
self-conjugate partitions of $n$.
\begin{enumerate}[label=(\alph*)]
\item Show that $p(n)\equiv sc(n)\pmod2$.
\item Show that $sc(n)$ is the number of partitions of $n$ into distinct odd parts, so that
$\sum_np(n)x^n\equiv\prod_{k\ge1}(1+x^{2k-1})\pmod2$.
\item Show that
\[
\prod_{k\ge1}(1+x^{2k-1})=\sum_{k\ge0}\frac{x^{k^2}}{(1-x^2)(1-x^4)\cdots(1-x^{2k})}.
\]
\end{enumerate}
\end{exercise}

\begin{exercise}[$\star\star$]\label{pb:glaisherthm}
\emph{Glaisher's theorem.} Let $k\ge2$.
\begin{enumerate}[label=(\alph*)]
\item Show that the partitions of $n$ with no part divisible by $k$ are as many as the partitions of $n$ in which no part
appears $k$ times or more.
\item Find an explicit bijection between the two sets, and apply it to the partition $7+2+2+2+2+1+1+1$ for $k=3$.
\end{enumerate}
\end{exercise}

\begin{exercise}[$\star\star$]\label{pb:eulerlog}
\emph{The logarithmic derivative.} Let $(a_d)_{d\ge1}$ be complex numbers and define $(b_n)$ by
\[
\sum_{n\ge0}b_nx^n=\prod_{d\ge1}(1-x^d)^{-a_d}.
\]
\begin{enumerate}[label=(\alph*)]
\item Show that $nb_n=\sum_{k=1}^nc_kb_{n-k}$ for $n\ge1$, where $c_k=\sum_{d\mid k}d\,a_d$.
\item Deduce that $np(n)=\sum_{k=1}^n\sigma(k)p(n-k)$, where $\sigma(k)$ is the sum of the divisors of $k$. Check it for $n=4$.
\item Let $M(d,q)$ be the number of monic irreducible polynomials of degree $d$ over $\mathbb F_q$. Using unique factorisation in
$\mathbb F_q[X]$, show that $\prod_d(1-x^d)^{-M(d,q)}=\frac1{1-qx}$, and deduce Gauss's formula $\sum_{d\mid k}d\,M(d,q)=q^k$.
\end{enumerate}
\end{exercise}

\begin{exercise}[$\star\star$]\label{pb:catalanodd}
\emph{Odd Catalan numbers.} Let $C_n=\frac1{n+1}\binom{2n}n$ be the Catalan numbers.
\begin{enumerate}[label=(\alph*)]
\item Show that $v_2(C_n)=s(n+1)-1$. Deduce that $C_n$ is odd if and only if $n=2^k-1$ for some $k\ge0$.
\item Give a second proof of the last statement from the equation $C(x)=1+xC(x)^2$ satisfied by $C(x)=\sum C_nx^n$, reduced
modulo $2$.
\end{enumerate}
\end{exercise}

\begin{exercise}[$\star\star$]\label{pb:stern}
\emph{Stern's diatomic sequence (1858).} Let $s(0)=0$, $s(1)=1$, $s(2n)=s(n)$ and $s(2n+1)=s(n)+s(n+1)$.
\begin{enumerate}[label=(\alph*)]
\item A \emph{hyperbinary representation} of $n$ writes $n=\sum_ie_i2^i$ with every $e_i\in\{0,1,2\}$. Show that $n$ has exactly
$s(n+1)$ of them.
\item Deduce that $\sum_{n\ge0}s(n+1)x^n=\prod_{k\ge0}\bigl(1+x^{2^k}+x^{2^{k+1}}\bigr)$.
\item Show that $\sum_{2^k\le n<2^{k+1}}s(n)=3^k$.
\item Show that $s(n)$ and $s(n+1)$ are coprime for every $n\ge0$.
\end{enumerate}
\end{exercise}

\begin{exercise}[$\star\star$]\label{pb:binarypart}
\emph{Binary partitions (Euler).} Let $b(n)$ be the number of partitions of $n$ into powers of $2$, so $b(0)=b(1)=1$, $b(2)=2$.
\begin{enumerate}[label=(\alph*)]
\item Show that $b(2n+1)=b(2n)$ and $b(2n)=b(2n-1)+b(n)$ for $n\ge1$.
\item Show in two ways that $b(n)$ is even for every $n\ge2$: from (a), and from the generating function, reduced modulo $2$.
\item Show that $b(2n)=b(0)+b(1)+\dots+b(n)$, and deduce that $\log b(n)\le\frac{(\log n)^2}{2\log2}+O(\log n)$.
\end{enumerate}
\end{exercise}

\begin{exercise}[$\star\star$]\label{pb:arcsin}
\emph{The square of the arcsine.}
\begin{enumerate}[label=(\alph*)]
\item Show that $y=\arcsin^2x$ satisfies $(1-x^2)y''-xy'=2$ on $]-1,1[$, with $y(0)=y'(0)=0$.
\item Deduce that for $|x|\le1$
\[
\arcsin^2x=\frac12\sum_{n\ge1}\frac{(2x)^{2n}}{n^2\binom{2n}n}.
\]
\item Deduce that $\displaystyle\sum_{n\ge1}\frac1{n^2\binom{2n}n}=\frac{\pi^2}{18}$ and $\displaystyle\sum_{n\ge1}\frac{4^n}{n^2\binom{2n}n}=\frac{\pi^2}2$.
\end{enumerate}
\end{exercise}

\begin{exercise}[$\star\star$]\label{pb:quotientlaw}
\emph{The law of $\{1/U\}$.} Let $U$ be uniform on $]0,1[$.
\begin{enumerate}[label=(\alph*)]
\item Show that $\{1/U\}$ has the density $g(u)=\sum_{k\ge1}\frac1{(k+u)^2}$ on $[0,1]$, which decreases from $\frac{\pi^2}6$ to $\frac{\pi^2}6-1$.
\item Show that $P\bigl(\lfloor1/U\rfloor\text{ is odd}\bigr)=\log2$ and $P\bigl(\{1/U\}<\frac12\bigr)=2-2\log2$.
\item Show that the proportion of $n\le x$ for which $\lfloor x/n\rfloor$ is odd tends to $\log2$ as $x\to\infty$.
\end{enumerate}
\end{exercise}

\begin{exercise}[$\star\star\star$]\label{pb:zigzag}
\emph{Up, down, up, down.} A permutation $\sigma$ of $\{1,\dots,n\}$ is \emph{alternating} if $\sigma(1)<\sigma(2)>\sigma(3)<\sigma(4)>\cdots$. Let $A_n$ be the number of
alternating permutations, with $A_0=A_1=1$; the values continue $1,2,5,16,61,272,1385$.
\begin{enumerate}[label=(\alph*)]
\item Show that $2A_{n+1}=\sum_{k=0}^n\binom nkA_kA_{n-k}$ for $n\ge1$. (Look at where $n+1$ sits, and count the permutations that zigzag starting either up or down.)
\item Let $f(z)=\sum_nA_n\frac{z^n}{n!}$. Show that $2f'=1+f^2$ and $f(0)=1$, and deduce that $f(z)=\sec z+\tan z$ for $|z|<\frac\pi2$.
\item Show that $f(z)-\frac4{\pi-2z}$ extends to a holomorphic function on the disc $|z|<\frac{3\pi}2$, and deduce that
\[
\frac{A_n}{n!}=2\Bigl(\frac2\pi\Bigr)^{n+1}\bigl(1+O(\delta^n)\bigr)\qq\text{for every }\delta>\tfrac13.
\]
\item What is the probability that a random permutation of $10$ elements is alternating? Compare with the approximation.
\end{enumerate}
\end{exercise}

\begin{exercise}[$\star\star$]\label{pb:poissonise}
\emph{Three people with the same birthday.} Draw uniformly and independently from $n$ values, and let $T$ be the number of draws needed
until some value has appeared $m$ times ($m\ge2$). Let $e_m(x)=\sum_{j<m}\frac{x^j}{j!}$.
\begin{enumerate}[label=(\alph*)]
\item Show that $P(T>k)=\frac{k!}{n^k}[x^k]\,e_m(x)^n$ for every $k\ge0$.
\item Deduce that $\mathbb E[T]=\int_0^\infty e^{-t}e_m(t/n)^n\dd t$.
\item For $m=2$, deduce that $\mathbb E[T]=1+Q(n)$, where $Q(n)=\sum_{k=1}^n\frac{n^{\underline k}}{n^k}$, and recover Volume~3, Exercise~2.7.
\item For $m=3$, show that $\mathbb E[T]\sim\Gamma(\frac43)(6n^2)^{1/3}$.
\item For $n=365$ the exact values are $24.62$ for a pair and $88.74$ for a triple. What does (d) give? Answer the think box.
\end{enumerate}
\end{exercise}

\begin{exercise}[$\star\star$]\label{pb:diehit}
\emph{Will the running total hit $n$?} Roll a fair die repeatedly and add the results. Let $u_n$ be the probability that some running total equals $n$, with $u_0=1$.
\begin{enumerate}[label=(\alph*)]
\item Show that $u_n=\frac16(u_{n-1}+\dots+u_{n-6})$ for $n\ge1$ (with $u_j=0$ for $j<0$), and that $\sum_nu_nz^n=\frac1{1-D(z)}$, where $D(z)=\frac16(z+z^2+\dots+z^6)$.
\item Show that $1-D$ has a simple root at $z=1$ and no other root in the closed unit disc.
\item Deduce that $u_n\to\frac27$, the inverse of the mean roll $\frac72$, and that $u_n-\frac27=O(\rho^n)$ for some $\rho<1$. (Numerically, one finds $\rho=0.7302$.)
\item Prove the same for any law on $\{1,\dots,N\}$ with mean $\mu$ whose support has greatest common divisor $1$: then $u_n\to\frac1\mu$. What goes wrong if the gcd is $2$?
\end{enumerate}
\end{exercise}

\begin{exercise}[$\star\star$]\label{pb:flory}
\emph{Dimers falling on a line.} A row has $n$ sites, numbered $1$ to $n$. Give each of the $n-1$ pairs $\{i,i+1\}$ an independent uniform label in $[0,1]$, and go
through the pairs in increasing order of labels, filling both sites of a pair whenever both are still empty. (At each step the pair filled is
uniform among the pairs still available.) Let $V_n$ be the expected number of sites left empty.
\begin{enumerate}[label=(\alph*)]
\item Show that $V_0=0$, $V_1=1$ and $(n-1)V_n=2\sum_{j=0}^{n-2}V_j$ for $n\ge2$.
\item Show that $\sum_{n\ge0}V_nx^n=\frac{xe^{-2x}}{(1-x)^2}$, so that $V_n=\sum_{k=0}^{n-1}(n-k)\frac{(-2)^k}{k!}$.
\item Deduce that $\bigl|V_n-(n+2)e^{-2}\bigr|\le\frac{2^{n+1}}{(n+1)!}$. In a long row, a fraction $e^{-2}=13.5\%$ of the sites stays empty.
\item Check the first-order result against the transfer theorem of Volume~4, Theorem~11.13.
\end{enumerate}
\end{exercise}

\section{Hints}

\begin{hint}{ex:binet}
$|\psi|<1$, so $\bigl|\frac{\psi^n}{\sqrt5}\bigr|<\frac12$. For the last identity, use $\varphi\psi=-1$.
\end{hint}
\begin{hint}{ex:change}
If $B=\frac A{1-x^c}$ then $B=A+x^cB$, that is, $b_n=a_n+b_{n-c}$.
\end{hint}
\begin{hint}{ex:dobinski}
$e^{e^x}=\sum_k\frac{e^{kx}}{k!}$, a double series with non-negative terms for $x\ge0$. For the recurrence, differentiate
$e^{e^x-1}$.
\end{hint}
\begin{hint}{ex:involutions}
Differentiate: $T'=(1+x)T$. Directly: in an involution of $\{1,\dots,n+1\}$, the element $n+1$ is fixed or swapped with one
of the others.
\end{hint}
\begin{hint}{ex:connected}
A graph is a set partition of its vertices into connected components, with a connected graph on each.
\end{hint}
\begin{hint}{ex:motzkin}
Look at the first step, and if it goes up, at the first return to the axis.
\end{hint}
\begin{hint}{ex:sigma}
$x\,D\log\prod(1-x^k)=-\sum_k\frac{kx^k}{1-x^k}=-\sum_{n\ge1}\sigma(n)x^n$. Multiply by $\prod(1-x^k)$ and use the theorem twice.
\end{hint}

\begin{hint}{pb:prouhet}
$\sum_{n<2^N}t_nx^n=\prod_{j<N}(1-x^{2^j})$ has a zero of order $N$ at $x=1$. Apply $(x\frac d{dx})^k$ and evaluate at $1$.
\end{hint}

\begin{hint}{pb:phipq}
(a) $\Phi_{pq}(x)=\frac{(1-x)(1-x^{pq})}{(1-x^p)(1-x^q)}$, and $\deg\Phi_{pq}<pq$: compute modulo $x^{pq}$. How many ways
can $m<pq$ be written as $ip+jq$ with $i,j\ge0$? (b) Use $\Phi_n=\prod_{d\mid n}(1-x^d)^{\mu(n/d)}$ modulo $x^8$.
\end{hint}

\begin{hint}{pb:commuting}
A $k$-tuple of commuting permutations is an action of $\Z^k$. Use the exponential formula (Volume~4, Chapter~2) on transitive
actions, and count the subgroups of index $m$ of $\Z^k$ by induction on $k$.
\end{hint}

\begin{hint}{pb:roots}
(a) Exponential formula. (b) Differentiate, and look at $(n-1)!/(n-p)!$ modulo $p$. (c) Elements of order $p$ come in groups
of $p-1$.
\end{hint}

\begin{hint}{pb:squares}
(a) What does squaring do to a cycle of odd length? Of even length? (b) Build the generating function cycle length by cycle
length: a factor $\exp$ for a free number of cycles, $\cosh$ for an even number.
\end{hint}

\begin{hint}{pb:randpoly}
(a) $f=g^2h$ with $h$ squarefree, in generating functions. (c) Choose $m_j$ irreducibles of degree $j$ with repetition.
\end{hint}

\begin{hint}{pb:parts123}
(a) $(1-x^6)=(1-x^3)(1+x^3)$, and $\frac1{(1-x)(1-x^2)}$ counts partitions into $1$'s and $2$'s. (b) Both sides change by $n$
when $n$ increases by $6$. (c) For a triangle with sides $x,y,z$, look at $y+z-x$, $x+z-y$, $x+y-z$.
\end{hint}

\begin{hint}{pb:selfconj}
(a) Conjugation is an involution. (b) Peel a self-conjugate diagram into hooks around the diagonal. (c) Classify by the size of
the largest square inside the diagram.
\end{hint}

\begin{hint}{pb:glaisherthm}
(a) $\frac{1-x^{kj}}{1-x^j}=1+x^j+\dots+x^{(k-1)j}$. (b) Write the multiplicity of each part in base $k$.
\end{hint}

\begin{hint}{pb:eulerlog}
(a) Apply $x\frac d{dx}\log$ to both sides. (c) A monic polynomial is a product of monic irreducibles in exactly one way.
\end{hint}

\begin{hint}{pb:catalanodd}
(a) $s(n+1)=s(n)+1-v_2(n+1)$. (b) Modulo $2$, $C(x)^2=C(x^2)$.
\end{hint}

\begin{hint}{pb:stern}
(a) Look at $e_0$. (c) Induction on $k$. (d) Euclid.
\end{hint}

\begin{hint}{pb:binarypart}
(a) Does the partition contain a $1$? (b) $(1-x)B(x)=B(x^2)$. (c) $b$ is non-decreasing.
\end{hint}

\begin{hint}{pb:arcsin}
(b) Find the recurrence for the Taylor coefficients, and for uniqueness note that $(\sqrt{1-x^2}\,D')'=0$ for the difference $D$.
(c) $x=\frac12$ and $x=1$.
\end{hint}

\begin{hint}{pb:quotientlaw}
(a) $\{1/U\}\le u$ and $\lfloor 1/U\rfloor=k$ mean $\frac1{k+u}\le U\le\frac1k$. (c) Count the $n$ with $\lfloor x/n\rfloor=k$, for $k\le K$.
\end{hint}

\begin{hint}{pb:zigzag}
(a) Removing $n+1$ from a permutation that zigzags (up first or down first) leaves two blocks; read the left one backwards. (b) Both sides of $2f'=1+f^2$ have the
same coefficients by (a); check that $\sec+\tan=\frac{1+\sin z}{\cos z}$ solves the equation, and compare coefficients. (c) The zeros of $\cos$ nearest to $0$ are $\pm\frac\pi2$; at
$-\frac\pi2$ the numerator $1+\sin z$ vanishes too. Use Cauchy's estimates (Volume~4, Corollary~9.12).
\end{hint}

\begin{hint}{pb:poissonise}
(a) Exponential generating functions count words. (b) $k!=\int_0^\infty t^ke^{-t}\dd t$, and the sum over $k$ is finite. (c) Expand $(1+\frac tn)^n$. (d) Substitute $t=nx$; the
function $e^{-x}(1+x+\frac{x^2}2)$ is $1-\frac{x^3}6+O(x^4)$ near $0$.
\end{hint}

\begin{hint}{pb:diehit}
(a) Condition on the first roll. (b) On $|z|\le1$, $|D(z)|\le1$, with equality in $D(z)=1$ only if all the $z^j$ are equal to $1$. (c) Subtract the pole at $1$, with residue computed from $D'(1)=\frac72$,
and use Cauchy's estimates on the rest.
\end{hint}

\begin{hint}{pb:flory}
(a) The pair with the smallest label is filled first and splits the row into two rows, whose labels are still independent and uniform given the first one. (b) Turn the
recurrence into $xV'-V=\frac{2x^2V}{1-x}$, and put $V=xW$. (c) The tail $\sum_{k\ge n}(k-n)\frac{(-2)^k}{k!}$ is alternating.
\end{hint}

\section{Solutions}

\begin{solution}{ex:binet}
\idea The term $\psi^n$ is a small correction that decays geometrically.
\writeup Since $|\psi|=\frac{\sqrt5-1}2<1$, $\bigl|F_n-\frac{\varphi^n}{\sqrt5}\bigr|=\frac{|\psi|^n}{\sqrt5}\le\frac1{\sqrt5}<\frac12$, so $F_n$ is the
integer nearest to $\frac{\varphi^n}{\sqrt5}$. Next,
$\frac{F_{n+1}}{F_n}=\frac{\varphi^{n+1}-\psi^{n+1}}{\varphi^n-\psi^n}=\varphi\cdot\frac{1-(\psi/\varphi)^{n+1}}{1-(\psi/\varphi)^n}\to\varphi$ since
$|\psi/\varphi|<1$. Finally, with $\varphi\psi=-1$ and $\varphi^2+\psi^2=(\varphi+\psi)^2-2\varphi\psi=1+2=3$,
\[
\begin{aligned}5F_{n-1}F_{n+1}&=\varphi^{2n}+\psi^{2n}-(\varphi\psi)^{n-1}(\varphi^2+\psi^2)\eqbr=\varphi^{2n}+\psi^{2n}+3(-1)^n,\\
5F_n^2&=\varphi^{2n}+\psi^{2n}-2(\varphi\psi)^n\eqbr=\varphi^{2n}+\psi^{2n}-2(-1)^n,\end{aligned}
\]
and subtracting, $5(F_{n-1}F_{n+1}-F_n^2)=5(-1)^n$. This is Cassini's identity.
\end{solution}

\begin{solution}{ex:change}
\idea The generating function is a product of one factor $\frac1{1-x^c}$ per coin, and each factor is cheap to apply.
\writeup Let $V=\{1,2,5,10,20,50,100\}$. The number of ways is $[x^{100}]\prod_{c\in V}\frac1{1-x^c}$, by the argument of
Proposition~\ref{prop:eulerproduct}: a choice of $m_c$ coins of value $c$ with $\sum cm_c=100$. If $B=\frac A{1-x^c}$ then
$(1-x^c)B=A$, that is, $b_n=a_n+b_{n-c}$ (with $b_n=0$ for $n<0$): a running sum with step $c$. Starting from
$a=(1,0,0,\dots)$ and applying the seven factors in turn, each in $101$ additions, gives $4563$. This is the classical
dynamic programming algorithm for the coin-change problem, and the generating function explains why the order of the
coins does not matter.
\end{solution}

\begin{solution}{ex:dobinski}
\idea Expand $e^{e^x}$ by the exponential series in the outer variable, then each $e^{kx}$ in $x$.
\writeup For real $x$, $e^{e^x}=\sum_k\frac{e^{kx}}{k!}=\sum_k\sum_n\frac{k^nx^n}{k!\,n!}$. For $x\ge0$ the terms are non-negative,
so the order of summation can be exchanged by Lemma~\ref{lem:packets}; for $x<0$ the double series is summable, being
dominated by the case $|x|$, and Lemma~\ref{lem:packets}(c) applies. Hence $e^{e^x-1}=\sum_n\bigl(\frac1e\sum_k\frac{k^n}{k!}\bigr)\frac{x^n}{n!}$ for every
real $x$, a power series converging on $\R$. By the remark at the end of the first section its coefficients are those of the
formal series $e^{e^x-1}=\sum b_n\frac{x^n}{n!}$, which gives Dobi\'nski's formula. For the recurrence,
$D\,e^{e^x-1}=e^x\,e^{e^x-1}$; the coefficient of $\frac{x^n}{n!}$ on the left is $b_{n+1}$ and on the right
$\sum_k\binom nkb_k\cdot1$, by the product rule for exponential series.
\end{solution}

\begin{solution}{ex:involutions}
\idea A differential equation for the generating function gives a recurrence for the coefficients.
\writeup With $T=e^{x+x^2/2}$, the chain rule gives $T'=(1+x)T$. The coefficient of $\frac{x^n}{n!}$ in $T'$ is $t_{n+1}$; in $T$ it
is $t_n$, and in $xT=\sum t_m\frac{x^{m+1}}{m!}$ it is $n\,t_{n-1}$. So $t_{n+1}=t_n+nt_{n-1}$. Directly: in an involution of
$\{1,\dots,n+1\}$, either $n+1$ is fixed, and the rest is any involution of $\{1,\dots,n\}$, or $n+1$ is swapped with one of
the $n$ others, $i$, and the rest is an involution of the remaining $n-1$ elements. The values are $1,1,2,4,10,26,76,232$.
\end{solution}

\begin{solution}{ex:connected}
\idea All graphs are assemblies of connected ones, so the exponential formula applies and its inverse is the logarithm.
\writeup A graph on $\{1,\dots,n\}$ is the same as the partition of the vertex set into its connected components, with a
connected graph on each block. By Theorem~\ref{thm:expformula}, $\sum_n2^{\binom n2}\frac{x^n}{n!}=\exp(\hat C)$, and by
Proposition~\ref{prop:explog} (third identity, with $1+F$ the series of all graphs), $\hat C=\log\sum_n2^{\binom n2}\frac{x^n}{n!}$.
Expanding $\log(1+F)=F-\frac{F^2}2+\frac{F^3}3-\cdots$ with $F=x+\frac{2x^2}2+\frac{8x^3}6+\frac{64x^4}{24}+\frac{1024x^5}{120}$ to order $5$ gives
\[
\hat C=x+\frac{x^2}2+\frac{2x^3}3+\frac{19x^4}{12}+\frac{91x^5}{15}+\cdots,
\]
so $c_1,\dots,c_5=1,1,4,38,728$. Check $c_3=4$: the path in three ways, and the triangle. Almost all graphs are connected:
$\frac{c_n}{2^{\binom n2}}\to1$.
\end{solution}

\begin{solution}{ex:motzkin}
\idea Decompose a path at its first step, and an up step at its matching down step, exactly as for Catalan numbers.
\writeup A non-empty Motzkin path either starts with a flat step followed by any Motzkin path, or starts with an up step;
in the second case, let the first return to the axis be at step $k+2$: between the up step and its matching down step
lies a Motzkin path of length $k$ (raised by one), and after it any Motzkin path. So $M=1+xM+x^2M^2$. Solving the quadratic
$x^2M^2+(x-1)M+1=0$ gives $M=\frac{1-x\pm\sqrt{(1-x)^2-4x^2}}{2x^2}$, and only the minus sign gives a power series (the plus
sign has a pole of order $2$ at $0$). Here $\sqrt{1-2x-3x^2}$ means the formal series $\exp\bigl(\frac12\log(1-2x-3x^2)\bigr)$,
whose square is $1-2x-3x^2$ by Proposition~\ref{prop:explog}. The recurrence $M_n=M_{n-1}+\sum_{k=0}^{n-2}M_kM_{n-2-k}$, read off
$M=1+xM+x^2M^2$, gives $1,1,2,4,9,21,51,127,323$.
\end{solution}

\begin{solution}{ex:sigma}
\idea A logarithmic derivative turns a product into a sum; for $\prod(1-x^k)$ the sum is a divisor sum, and for the
pentagonal series it is a sparse sum.
\writeup Let $E=\prod(1-x^k)$. The logarithm of an infinite product of series $1+A_k$ (with orders of $A_k$ tending to
infinity) is the sum of the logarithms, by Proposition~\ref{prop:explog} applied to finite products and comparison of
coefficients. So $\log E=\sum_k\log(1-x^k)$ and
\[
\begin{aligned}x\frac{DE}E&=\sum_kx\frac{-kx^{k-1}}{1-x^k}\eqbr=-\sum_k\sum_{m\ge1}kx^{km}\eqbr=-\sum_{n\ge1}\sigma(n)x^n,\end{aligned}
\]
grouping the terms with $km=n$. Hence $-E\sum\sigma(n)x^n=xDE$. By Theorem~\ref{thm:pentagonal}, $E=\sum_j(-1)^jx^{g_j}$ and
$xDE=\sum_j(-1)^jg_jx^{g_j}$, with $g_j=\frac{j(3j-1)}2$. Comparing coefficients of $x^n$:
\[
-\sum_j(-1)^j\sigma(n-g_j)=(-1)^{j_0}n\ \text{ if }n=g_{j_0},\ \text{else }0,
\]
with $\sigma(m)=0$ for $m\le0$. The term $j=0$ is $\sigma(n)$, so $\sigma(n)=\sum_{j\ne0}(-1)^{j+1}\sigma(n-g_j)$, plus $-(-1)^{j_0}n$
when $n=g_{j_0}$; this is the convention ``$\sigma(0)=n$'' in the term $-(-1)^{j_0}\sigma(n-g_{j_0})$. For instance $\sigma(5)=\sigma(4)+\sigma(3)-5=7+4-5=6$.
\end{solution}

\begin{solution}{pb:prouhet}
\idea The signs $t_n$ are built digit by digit, so their generating function is a product with one factor per binary
digit, and each factor vanishes at $x=1$. Moments $\sum t_nn^k$ are derivatives at $1$. A zero of order $N$ kills all
derivatives of order below $N$.

\writeup Each $n<2^N$ is $\sum_{j<N}\varepsilon_j2^j$ with $\varepsilon_j\in\{0,1\}$, and $t_n=\prod_j(-1)^{\varepsilon_j}$. Expanding
the product, $G(x)=\sum_{n<2^N}t_nx^n=\prod_{j=0}^{N-1}(1-x^{2^j})$. Each factor vanishes at $1$, so $G=(1-x)^NH(x)$ with
$H(1)=\prod_j(1+1+\dots+1)=\prod_j2^j=2^{N(N-1)/2}$, since $\frac{1-x^{2^j}}{1-x}=1+x+\dots+x^{2^j-1}$.

Let $D=x\frac d{dx}$. Then $D^kG(x)=\sum_nt_nn^kx^n$, so the sum we want is $(D^kG)(1)$. By induction on $k$,
$D^kG=\sum_{i=0}^ka_{k,i}x^iG^{(i)}$ with integers $a_{k,i}$ and $a_{k,k}=1$: apply $D$ to $x^iG^{(i)}$ to get
$ix^iG^{(i)}+x^{i+1}G^{(i+1)}$. Since $G$ has a zero of order $N$ at $1$, $G^{(i)}(1)=0$ for $i<N$.

(a) For $k<N$, every term of $(D^kG)(1)$ contains some $G^{(i)}(1)$ with $i\le k<N$, so it is $0$.

(b) For $k=N$, only the term $i=N$ survives: the sum is $G^{(N)}(1)$. From $G=(-1)^N(x-1)^NH$ and Leibniz's rule,
$G^{(N)}(1)=(-1)^NN!\,H(1)=(-1)^NN!\,2^{N(N-1)/2}$.

(c) Take $N=4$. The numbers $n<16$ with $t_n=1$ (an even number of binary ones) are
$A=\{0,3,5,6,9,10,12,15\}$, and the others are $B=\{1,2,4,7,8,11,13,14\}$. By (a), $\sum_Ad^k=\sum_Bd^k$ for $k=0,1,2,3$:
the common values are $8$, $60$, $620$ and $7200$.
\tools Volume~3, Chapter~3 and Volume~4, Chapter~2 (generating functions); Volume~0, Chapter~3 (multiplicity of roots).
\leads This is Prouhet's solution (1851) of what is now called the Prouhet--Tarry--Escott problem: find two different sets
with equal power sums up to degree $k$. The Thue--Morse sequence uses $2^{k}$ numbers in each set; the best possible size is
$k+1$ numbers in each set, and such \emph{ideal} solutions are known for $k\le9$ and $k=11$, but not for $k=10$ or any $k\ge12$. The same sequence appears in fair sharing (``who chooses next''),
in the Rudin--Shapiro-like signs of Chapter~1, and in symbolic dynamics.
\end{solution}

\begin{solution}{pb:phipq}
\idea $\Phi_{pq}$ is a ratio of four binomials $1-x^d$. The denominator $\frac1{(1-x^p)(1-x^q)}$ is the generating function
of the numbers $ip+jq$, and below $pq$ each number has at most one such representation. Since $\Phi_{pq}$ has degree less
than $pq$, we lose nothing by working modulo $x^{pq}$, where $1-x^{pq}$ is $1$. For $\Phi_{105}$ the same reduction modulo
$x^8$ leaves a two-line computation.

\writeup From $x^n-1=\prod_{d\mid n}\Phi_d$ and Möbius inversion in multiplicative form,
$\Phi_n=\prod_{d\mid n}(x^d-1)^{\mu(n/d)}$, and for $n>1$, $\sum_{d\mid n}\mu(n/d)=0$, so the signs can be changed:
$\Phi_n=\prod_{d\mid n}(1-x^d)^{\mu(n/d)}$, an identity of rational functions, and of power series in $\Z[[x]]$.

(a) For $n=pq$: $\Phi_{pq}=\frac{(1-x)(1-x^{pq})}{(1-x^p)(1-x^q)}$. In $\Z[[x]]$,
$\frac1{(1-x^p)(1-x^q)}=\sum_{i,j\ge0}x^{ip+jq}$. If $ip+jq=i'p+j'q<pq$ then $p(i-i')=q(j'-j)$, so $q\mid i-i'$; as
$0\le i,i'<q$, $i=i'$ and then $j=j'$. So modulo $x^{pq}$,
\[
\Phi_{pq}\equiv(1-x)\sum_{m<pq}\mathbf 1_S(m)x^m,\qquad S=\{ip+jq:i,j\ge0\}.
\]
Since $\deg\Phi_{pq}=(p-1)(q-1)<pq$, the coefficient of $x^m$ in $\Phi_{pq}$ is exactly $\mathbf 1_S(m)-\mathbf 1_S(m-1)$,
for every $m<pq$. This is $-1$, $0$ or $1$. Moreover it is $1$ where $\mathbf 1_S$ jumps up and $-1$ where it jumps down,
and up and down jumps alternate: the non-zero coefficients alternate between $1$ and $-1$.

(b) $n=105=3\cdot5\cdot7$. Modulo $x^8$ only the divisors $d<8$ matter: $d=1,3,5,7$, with $\mu(105)=-1$, $\mu(35)=\mu(21)=\mu(15)
=1$. So
\[
\Phi_{105}\equiv\frac{(1-x^3)(1-x^5)(1-x^7)}{1-x}\equiv(1+x+\dots+x^7)(1-x^3-x^5-x^7)\pmod{x^8},
\]
since $x^8\equiv0$. The coefficient of $x^7$ is $1-1-1-1=-2$.
\tools Volume~1, Chapters~4 and~5 (Möbius inversion, cyclotomic polynomials); Volume~3, Chapter~3 and Volume~4, Chapter~2
(generating functions, the coin problem).
\leads The set $S$ is the set of amounts payable with coins $p$ and $q$; its complement has $\frac{(p-1)(q-1)}2$ elements, the
largest being $pq-p-q$ (Sylvester's coin problem, 1884), and this is the degree of $\Phi_{pq}$ minus the last jump. This
proof is due to Lam and Leung (1996). Schur's proof that the coefficients are unbounded uses $\Phi_{p_1\cdots p_t}$ with
$t$ odd and $p_1<\dots<p_t<p_1+p_2$: its coefficient of $x^{p_t}$ is $1-t$.
\end{solution}

\begin{solution}{pb:commuting}
\idea Commuting permutations $\sigma_1,\dots,\sigma_k$ are the same thing as an action of $\Z^k$ on $\{1,\dots,n\}$, the $i$-th
generator acting by $\sigma_i$. An action splits into transitive pieces, and the exponential formula turns ``a set of
transitive pieces'' into an exponential. A transitive action of an abelian group on $m$ points is a subgroup of index $m$ plus a
labelling. So everything reduces to counting subgroups of $\Z^k$ of index $m$.

\writeup Let $\Gamma=\Z^k$. Commuting $k$-tuples of permutations of a set $S$ correspond to actions of $\Gamma$ on $S$. An action
is the disjoint union of its restrictions to its orbits, which are transitive actions. Let $t_m$ be the number of transitive
actions of $\Gamma$ on $\{1,\dots,m\}$. By the exponential formula (Volume~4, Chapter~2),
\[
\sum_{n\ge0}h_k(n)\frac{x^n}{n!}=\exp\Bigl(\sum_{m\ge1}t_m\frac{x^m}{m!}\Bigr).
\]
\emph{Transitive actions.} Given a transitive action on $\{1,\dots,m\}$, let $H$ be the stabiliser of the point $1$: a subgroup
of index $m$ (orbit--stabiliser), and $\gamma H\mapsto\gamma\cdot1$ is a bijection $\Gamma/H\to\{1,\dots,m\}$ sending $H$ to $1$.
Conversely a subgroup $H$ of index $m$ and a bijection $\beta:\Gamma/H\to\{1,\dots,m\}$ with $\beta(H)=1$ define a transitive action
$\gamma\cdot\beta(\delta H)=\beta(\gamma\delta H)$, whose stabiliser of $1$ is $H$. These are inverse constructions, and there are
$(m-1)!$ bijections $\beta$. So $t_m=(m-1)!\,a_m(\Gamma)$, with $a_m$ the number of subgroups of index $m$, and
\[
\sum_nh_k(n)\frac{x^n}{n!}=\exp\Bigl(\sum_{m\ge1}a_m(\Z^k)\frac{x^m}m\Bigr).
\]
\emph{Subgroups of $\Z^k$ of index $m$.} We show $a_m(\Z^k)=\sum_{e\mid m}e\,a_e(\Z^{k-1})$. Let $H\le\Z^k$ have index $m$, write
$\Z^k=\Z^{k-1}\times\Z$, let $f\Z$ be the image of $H$ under the last coordinate ($f\ge1$ since the index is finite) and
$H'=H\cap(\Z^{k-1}\times\{0\})$. Then $H=H'\oplus\Z v$ for any $v=(w,f)\in H$, the index of $H'$ in $\Z^{k-1}$ is $e=\frac mf$, and $w$ is
determined modulo $H'$ by $H$. Conversely, for $f\mid m$, $H'$ of index $\frac mf$ in $\Z^{k-1}$ and any $w$, the group
$H'\oplus\Z(w,f)$ has last-coordinate image $f\Z$, meets $\Z^{k-1}\times\{0\}$ in $H'$ (an element $h'+t(w,f)$ has last coordinate
$tf$), and has index $\frac mf\cdot f=m$; and two choices $w,w'$ give the same group if and only if $w-w'\in H'$, which leaves
$e=\frac mf$ choices. Hence the recursion. With $a_m(\Z)=1$: $a_m(\Z^2)=\sigma(m)$ and $a_m(\Z^3)=\sum_{e\mid m}e\,\sigma(e)$.

\emph{The products.} For $k=2$, writing $m=de$,
\[
\sum_m\sigma(m)\frac{x^m}m=\sum_{d,e\ge1}d\,\frac{x^{de}}{de}=\sum_{e\ge1}-\log(1-x^e),
\]
so the exponential is $\prod_e(1-x^e)^{-1}$. For $k=3$, writing $m=de$,
\[
\sum_m\Bigl(\sum_{d\mid m}d\,\sigma(d)\Bigr)\frac{x^m}m=\sum_{d,e}\sigma(d)\frac{x^{de}}e=\sum_d-\sigma(d)\log(1-x^d),
\]
so the exponential is $\prod_d(1-x^d)^{-\sigma(d)}$. The first product is $\sum p(n)x^n$ (Volume~4, Chapter~2), so $h_2(n)=n!p(n)$.
\tools Volume~2, Chapters~3 and~7 (orbit--stabiliser, subgroups of $\Z^k$); Volume~4, Chapter~2 (exponential formula, partitions); Volume~1, Chapter~4 (divisor sums).
\leads The same argument gives, for every $k$,
\[
\sum_nh_k(n)\frac{x^n}{n!}=\exp\Bigl(\sum_ma_m(\Z^k)\frac{x^m}m\Bigr),\qquad\sum_m\frac{a_m(\Z^k)}{m^s}=\zeta(s)\zeta(s-1)\cdots\zeta(s-k+1).
\] Bryan and Fulman (1998) found these products as ``orbifold Euler
characteristics'' in string theory. $h_2(n)=n!p(n)$ also follows from Volume~2, Exercise~3.8: the number of
commuting pairs is $|G|$ times the number of classes, and $\mathfrak S_n$ has $p(n)$ classes.
\end{solution}

\begin{solution}{pb:roots}
\idea $\sigma^k=e$ means that every cycle length divides $k$, a restriction on the cycle lengths, which the exponential formula
handles. Differentiating gives a recursion, and modulo $p$ the coefficient $\frac{(n-1)!}{(n-p)!}$, a product of $p-1$
consecutive integers, is $0$ or $-1$.

\writeup (a) A permutation is a set of cycles, and there are $(d-1)!$ cycles on a $d$-set. By the exponential formula
(Volume~4, Chapter~2), permutations whose cycle lengths lie in a set $D$ have exponential generating function
$\exp\bigl(\sum_{d\in D}(d-1)!\frac{x^d}{d!}\bigr)=\exp\bigl(\sum_{d\in D}\frac{x^d}d\bigr)$. Take $D$ the divisors of $k$.

(b) For $k=p$, $A(x)=\exp(x+\frac{x^p}p)$ and $A'=(1+x^{p-1})A$. Comparing coefficients of $\frac{x^{n-1}}{(n-1)!}$:
$a_n=a_{n-1}+\frac{(n-1)!}{(n-p)!}a_{n-p}$. The product $(n-p+1)\cdots(n-1)$ runs over $p-1$ consecutive integers, which cover all
residues modulo $p$ except that of $n$. If $p\nmid n$, one of them is divisible by $p$ and the product is $\equiv0$; if $p\mid n$,
they are the non-zero residues and the product is $\equiv(p-1)!\equiv-1$ (Wilson). So $a_n\equiv a_{n-1}$ if $p\nmid n$ and
$a_n\equiv a_{n-1}-a_{n-p}$ if $p\mid n$. Since $a_0=\dots=a_{p-1}=1$, $a_p\equiv1-1=0$, and by induction all later terms vanish
modulo $p$: if $p\nmid n$, $a_n\equiv a_{n-1}\equiv0$; if $p\mid n$ and $n\ge2p$, $a_n\equiv0-0$.

(c) The solutions of $\sigma^p=e$ other than $e$ have order $p$. Each generates a cyclic group of order $p$, which contains
$p-1$ elements of order $p$, and two such groups are equal or meet in $\{e\}$. So $a_n-1$ is a multiple of $p-1$. With
$p\mid a_n$ and $\gcd(p,p-1)=1$, $a_n\equiv p\pmod{p(p-1)}$, since $p$ satisfies both congruences.
\tools Volume~4, Chapter~2 (exponential formula); Volume~1, Chapter~3 (Wilson); Volume~2, Chapter~1 (cyclic subgroups).
\leads Frobenius (1895) proved that in every finite group the number of solutions of $x^m=e$ is a multiple of
$\gcd(m,|G|)$; (b) is the case $\mathfrak S_n$, $m=p$. A famous conjecture of Frobenius, proved in 1991 with the
classification of finite simple groups, says that if there are exactly $m$ solutions, they form a subgroup.
\end{solution}

\begin{solution}{pb:squares}
\idea Squaring keeps an odd cycle as one cycle of the same length and breaks an even cycle into two equal halves. So even
cycles of a square come in pairs. That is a condition on the cycle type, cycle length by cycle length, and the exponential
formula still works if, for even lengths, we allow only even numbers of cycles.

\writeup (a) Let $\tau$ have a cycle $(a_1\,a_2\,\dots\,a_\ell)$. If $\ell$ is odd, $\tau^2$ maps $a_i\mapsto a_{i+2}$, and since $2$ is
invertible modulo $\ell$ this is a single $\ell$-cycle. If $\ell=2j$, $\tau^2$ splits it into $(a_1\,a_3\,\dots)$ and
$(a_2\,a_4\,\dots)$, two $j$-cycles. So in $\sigma=\tau^2$, the cycles of an even length $j$ come only from $2j$-cycles of $\tau$,
two at a time: their number is even. Conversely, suppose the condition holds. An odd cycle $c$ of length $\ell$ is a square:
$c=(c^{(\ell+1)/2})^2$. Two cycles $(a_1\,\dots\,a_j)$ and $(b_1\,\dots\,b_j)$ of the same length are the square of the
$2j$-cycle $(a_1\,b_1\,a_2\,b_2\,\dots\,a_j\,b_j)$. Group the even cycles of $\sigma$ in pairs of equal length, take these roots,
and multiply them (they have disjoint supports, so they commute).

(b) A permutation is determined by its set of cycles, and by the exponential formula with cycle lengths treated separately,
the generating function of permutations with $m_\ell$ cycles of length $\ell$ allowed only for $m_\ell$ in a set $S_\ell$ is
$\prod_\ell\sum_{m\in S_\ell}\frac1{m!}\bigl(\frac{x^\ell}\ell\bigr)^m$. (The coefficient of $\frac{x^n}{n!}$ counts permutations with
those cycle types, by Cauchy's formula $\frac{n!}{\prod\ell^{m_\ell}m_\ell!}$.) For odd $\ell$ every $m$ is allowed, giving
$\exp\frac{x^\ell}\ell$; for even $\ell$ only even $m$, giving $\cosh\frac{x^\ell}\ell$. The odd factors multiply to $\sqrt{\frac{1+x}{1-x}}$
as in Problem~\ref{pb:oddcycles}.

(c) $\sqrt{\frac{1+x}{1-x}}=\frac{1+x}{\sqrt{1-x^2}}$, and the other factors are even functions of $x$. So the generating
function is $(1+x)E(x)$ with $E$ even, and the coefficients of $x^{2m}$ and $x^{2m+1}$ are equal: $\frac{s_{2m+1}}{(2m+1)!}=\frac{s_{2m}}{(2m)!}$.
And a permutation with all cycles odd is a square by (a), so $s_n\ge o_n$.
\tools Volume~2, Chapter~4 (cycles, powers); Volume~4, Chapter~2 (exponential formula).
\leads Blum (1974) proved that the probability is $\sim\sqrt{\frac2{\pi n}}\prod_{m\ge1}\cosh\frac1{2m}=1.2218\sqrt{\frac2{\pi n}}$. The
a direct computation verifies this up to $n=800$, where the ratio is $1.002$. Pouyanne (2002) treated $k$-th roots.
\end{solution}

\begin{solution}{pb:randpoly}
\idea Polynomials over $\mathbb F_q$ factor uniquely, so their counting functions satisfy Euler products, exactly like the
integers, but with power series in $t=q^{-s}$ that can be summed in closed form.

\writeup (a) Every monic $f$ is uniquely $g^2h$ with $g$, $h$ monic and $h$ squarefree (group the irreducible factors by parity of
their multiplicity). With $s_n$ the number of squarefree monic polynomials of degree $n$ and $S(t)=\sum s_nt^n$, this says
$\frac1{1-qt}=S(t)\cdot\frac1{1-qt^2}$, so $S(t)=\frac{1-qt^2}{1-qt}$ and $s_n=q^n-q^{n-1}$ for $n\ge2$.

(b) By Volume~2, Theorem~9.10, $nM(n,q)=\sum_{d\mid n}\mu(d)q^{n/d}$. The term $d=1$ is $q^n$, and the others have $n/d\le n/2$, so
$|nM(n,q)-q^n|\le\sum_{e\le n/2}q^e\le\frac q{q-1}q^{n/2}\le2q^{n/2}$. Divide by $nq^n$.

(c) Choosing $f$ of type $\lambda$ means choosing, for each $j$, a multiset of $m_j$ monic irreducibles of degree $j$: there are
$\prod_j\binom{M(j,q)+m_j-1}{m_j}$ such $f$. As $q\to\infty$, $M(j,q)\sim\frac{q^j}j$ by (b), so
$\binom{M+m-1}m\sim\frac{M^m}{m!}\sim\frac{q^{jm}}{j^mm!}$, and the product is $\sim q^{\sum jm_j}/\prod_jj^{m_j}m_j!=q^n/z_\lambda$. And
$1/z_\lambda$ is the probability that a uniform permutation of $\mathfrak S_n$ has cycle type $\lambda$, because the class of type
$\lambda$ has $n!/z_\lambda$ elements (Volume~2, Chapter~4).

\tools Volume~2, Chapter~9 (counting irreducibles); Volume~4, Chapter~2 (generating functions); Volume~2, Chapter~4 (conjugacy classes of
$\mathfrak S_n$).
\leads Compare with the integers: squarefree density $1/\zeta(2)=6/\pi^2$, and here $1/Z(q^{-2})=1-\frac1q$ where
$Z(t)=\frac1{1-qt}$ plays the role of $\zeta$. The error $q^{-n/2}$ in (b) is the Riemann hypothesis for $\mathbb F_q[t]$, which is
easy; for the integers it is the open problem. The analogy ``factor degrees = cycle lengths'' is the Frobenius acting on the roots,
and Chebotarev's theorem makes it precise.
\end{solution}

\begin{solution}{pb:parts123}
\idea The generating function $\frac1{(1-x)(1-x^2)(1-x^3)}$ has a pole of order $3$ at $1$, which gives the quadratic
$\frac{(n+3)^2}{12}$, and smaller poles at roots of unity of order $2$, $3$, $6$, whose contributions are periodic and small. A
difference of step $6$ kills them all.

\writeup (a) The generating function of $a$ is $A(x)=\frac1{(1-x)(1-x^2)(1-x^3)}$, so
\[
(1-x^6)A(x)=\frac{1+x^3}{(1-x)(1-x^2)}.
\]
The coefficient of $x^n$ in $\frac1{(1-x)(1-x^2)}$ is the number of ways to write $n=u+2v$, which is $\lfloor\frac n2\rfloor+1$. So for
$n\ge6$,
\[
a(n)-a(n-6)=\Bigl\lfloor\frac n2\Bigr\rfloor+1+\Bigl\lfloor\frac{n-3}2\Bigr\rfloor+1=n,
\]
as one checks separately for $n$ even ($\frac n2+1+\frac{n-4}2+1$) and odd ($\frac{n-1}2+1+\frac{n-3}2+1$).

(b) Let $g(n)=\frac{(n+3)^2}{12}$. Then $g(n)-g(n-6)=\frac{(n+3)^2-(n-3)^2}{12}=n$, so $a(n)-g(n)$ depends only on $n\bmod6$. For
$n=0,\dots,5$ we have $a(n)=1,1,2,3,4,5$ and $g(n)=\frac34,\frac43,\frac{25}{12},3,\frac{49}{12},\frac{16}3$, so
$a(n)-g(n)\in\{\frac14,-\frac13,-\frac1{12},0,-\frac1{12},-\frac13\}$. These are all less than $\frac12$ in absolute value, so $a(n)$ is
the integer nearest to $g(n)$, for every $n$.

(c) Let $x\le y\le z$ be the sides, $x+y+z=n$, $x+y>z$. Put $\alpha=y+z-x\ge\beta=x+z-y\ge\gamma=x+y-z\ge1$. These are
positive, have the parity of $n$, add up to $n$, and determine $(x,y,z)=(\frac{\beta+\gamma}2,\frac{\alpha+\gamma}2,\frac{\alpha+\beta}2)$;
conversely any such triple gives a triangle. If $n=2m+3$, write $\alpha=2a'+1$, $\beta=2b'+1$, $\gamma=2c'+1$: then
$a'\ge b'\ge c'\ge0$ and $a'+b'+c'=m$, a partition of $m$ into at most $3$ parts. By conjugation these are as many as the
partitions of $m$ into parts at most $3$, which are counted by $a(m)$. If $n=2m$, write $\alpha=2a''$, $\beta=2b''$, $\gamma=2c''$
with $a''\ge b''\ge c''\ge1$ and sum $m$; subtracting $1$ from each gives a partition of $m-3$ into at most $3$ parts.

\tools Volume~4, Chapter~2 (generating functions, partial fractions, conjugation).
\leads The sequence of (c) is known as Alcuin's sequence: it also counts the solutions of a cask-sharing puzzle in the
\emph{Propositiones ad acuendos juvenes}, a collection attributed to Alcuin of York (OEIS A005044). For parts $1,\dots,k$ the
same method gives $a_k(n)=\frac{n^{k-1}}{k!(k-1)!}+O(n^{k-2})$, but the nearest-integer formula is special to small $k$.
\end{solution}

\begin{solution}{pb:selfconj}
\idea An involution changes the parity of a finite set only through its fixed points. The fixed points of conjugation have a
rigid shape that is easy to count.

\writeup (a) Conjugation is an involution of the set of partitions of $n$. The partitions that are not fixed come in pairs
$\{\lambda,\lambda'\}$, so $p(n)$ and the number $sc(n)$ of fixed points have the same parity.

(b) Let $\lambda$ be self-conjugate, and let $k$ be the number of dots on the diagonal of its Ferrers diagram. The $i$-th
\emph{diagonal hook} consists of the dot $(i,i)$, the dots to its right in row $i$, and the dots below it in column $i$. Since the
diagram is symmetric, the arm and the leg of this hook have the same length $\lambda_i-i$, so the hook has $2(\lambda_i-i)+1$ dots, an
odd number, and these numbers strictly decrease with $i$. The hooks partition the diagram, so we get a partition of $n$ into
distinct odd parts. Conversely, distinct odd parts $h_1>\dots>h_k$ can be bent into nested symmetric hooks, each fitting outside
the previous one because $h_i\ge h_{i+1}+2$; this is the inverse map. The generating function of partitions into distinct odd
parts is $\prod_k(1+x^{2k-1})$.

(c) A self-conjugate partition whose largest inscribed square (the Durfee square) is $k\times k$ consists of that square,
of the partition $\mu$ formed by the rows below the square, whose parts are at most $k$, and of its mirror image $\mu'$ to the
right of the square. So it is determined by $k$ and by a partition $\mu$ into parts at most $k$, and its size is $k^2+2|\mu|$. Summing over $k$,
the generating function of self-conjugate partitions is $\sum_kx^{k^2}\prod_{j\le k}(1-x^{2j})^{-1}$. Compare with (b).

\tools Volume~4, Chapter~2 (Ferrers diagrams, conjugation).
\leads Parity of $p(n)$ is a famous mystery: it is conjectured that $p(n)$ is odd for asymptotically half of all $n$, but it is
not even known that the odd values have a positive lower density. Problem~\ref{pb:kolberg} shows at least that both parities occur
infinitely often.
\end{solution}

\begin{solution}{pb:glaisherthm}
\idea Generating functions give the equality of the counts in one line. For the bijection, a part repeated $m$ times can be
merged into bigger parts according to the base-$k$ digits of $m$.

\writeup (a) The generating function of partitions with no part divisible by $k$ is $\prod_{k\nmid j}(1-x^j)^{-1}$; that of
partitions with every multiplicity at most $k-1$ is
\[
\prod_{j\ge1}(1+x^j+\dots+x^{(k-1)j})=\prod_{j\ge1}\frac{1-x^{kj}}{1-x^j}.
\]
In the last product, every factor $1-x^{kj}$ of the numerator cancels the factor of the denominator with index $kj$, and what is
left is $\prod_{k\nmid j}(1-x^j)^{-1}$. The two products are equal in $\Z[[x]]$, so the counts agree for every $n$.

(b) Start from a partition with no part divisible by $k$. If the part $j$ appears $m$ times, write $m=\sum_ia_ik^i$ in base $k$, and
replace the $m$ copies of $j$ by $a_i$ copies of $jk^i$ for each $i$; the sum is unchanged. Every integer is uniquely $jk^i$ with
$k\nmid j$, so the new parts $jk^i$ are distinct for distinct pairs $(j,i)$, and each appears $a_i\le k-1$ times. Conversely, from a
partition with multiplicities at most $k-1$, write each part as $jk^i$ with $k\nmid j$, and break each copy of $jk^i$ into $k^i$
copies of $j$; the base-$k$ expansion of the resulting multiplicity of $j$ recovers the digits. The two maps are inverse.

For $7+2+2+2+2+1+1+1$ (no part divisible by $3$): the part $7$ appears once, $1=1_3$, and stays; the part $2$ appears $4=11_3$
times and becomes $2+6$; the part $1$ appears $3=10_3$ times and becomes $3$. The image is $7+6+3+2$, a partition of $18$ with all
multiplicities at most $2$, as is the original.

\tools Volume~4, Chapter~2 (generating functions of partitions); Volume~4, Proposition~2.11 (the case $k=2$).
\leads For $k=2$ this is Euler's ``distinct parts equal odd parts'', and (b) is its binary bijection; Sylvester found a very
different one. Glaisher published the general case in 1883. Many partition identities of this kind, such as the Rogers--Ramanujan identities, have generating
function proofs long before any bijection was found.
\end{solution}

\begin{solution}{pb:eulerlog}
\idea The logarithm turns the infinite product into a sum, and $x\frac d{dx}$ turns it back into a power series with integer
structure. Multiplying back by $B(x)$ gives a recurrence.

\writeup (a) Let $B(x)=\sum b_nx^n$. As formal power series with constant term $1$,
\[
\log B(x)=\sum_da_d\log\frac1{1-x^d}=\sum_da_d\sum_{r\ge1}\frac{x^{dr}}r,\qquad x\frac{B'(x)}{B(x)}=\sum_da_d\sum_{r\ge1}dx^{dr}=\sum_{k\ge1}c_kx^k,
\]
with $c_k=\sum_{d\mid k}da_d$ (put $k=dr$). So $xB'(x)=B(x)\sum_kc_kx^k$, and comparing coefficients of $x^n$ gives
$nb_n=\sum_{k=1}^nc_kb_{n-k}$. (One can avoid logarithms: the logarithmic derivative of a product is the sum of the logarithmic
derivatives of the factors, and $x\frac d{dx}$ of $(1-x^d)^{-a_d}$ divided by itself is $\frac{a_ddx^d}{1-x^d}$.)

(b) Take $a_d=1$ for all $d$: then $b_n=p(n)$ and $c_k=\sigma(k)$. For $n=4$: $\sigma(1)p(3)+\sigma(2)p(2)+\sigma(3)p(1)+\sigma(4)p(0)
=3+3\cdot2+4+7=20=4\,p(4)$.

(c) Expand $\prod_P(1-x^{\deg P})^{-1}$ over all monic irreducible $P$: a term $x^n$ arises from each choice of finitely many $P$ with
multiplicities and total degree $n$, that is, by unique factorisation, from each monic polynomial of degree $n$. There are $q^n$ of
them, so the product is $\sum q^nx^n=\frac1{1-qx}$, and grouping the $P$ by degree gives $\prod_d(1-x^d)^{-M(d,q)}$. Now apply
(a) with $a_d=M(d,q)$, $b_n=q^n$; or directly: $x\frac{d}{dx}\log\frac1{1-qx}=\frac{qx}{1-qx}=\sum_kq^kx^k$, so $c_k=q^k$, that is
$\sum_{d\mid k}dM(d,q)=q^k$.

\tools Volume~4, Chapter~2 (formal logarithm and derivative); Volume~2, Chapter~9 (irreducible polynomials over $\mathbb F_q$).
\leads Möbius inversion turns Gauss's formula into $M(k,q)=\frac1k\sum_{d\mid k}\mu(k/d)q^d$. The map $(a_d)\mapsto(b_n)$ is the
\emph{Euler transform}; it turns ``connected'' objects into all objects whenever objects factor uniquely into connected ones,
as integers factor into primes and polynomials into irreducibles.
\end{solution}

\begin{solution}{pb:catalanodd}
\idea Kummer's theorem turns the parity question into a statement about binary digits; the functional equation gives the same
answer through the Frobenius map $C\mapsto C^2$.

\writeup (a) By Problem~\ref{pb:kummer}, $v_2\binom{2n}n$ is the number of carries in $n+n$ in binary, which is $s(n)$ (each digit $1$ of
$n$ produces one carry, since doubling is a shift). So $v_2(C_n)=s(n)-v_2(n+1)$. Adding $1$ to $n$ turns the final block of
$v_2(n+1)$ ones into zeros and the next zero into a one, so $s(n+1)=s(n)-v_2(n+1)+1$. Hence $v_2(C_n)=s(n+1)-1$, which vanishes
exactly when $n+1$ is a power of $2$.

(b) Reduce $C=1+xC^2$ modulo $2$. In $\mathbb F_2[[x]]$, squaring is additive, so $C(x)^2=\sum C_nx^{2n}=C(x^2)$, and
$C(x)\equiv1+xC(x^2)$. In coefficients: $C_0\equiv1$, $C_{2m}\equiv0$ for $m\ge1$, and $C_{2m+1}\equiv C_m$. By induction $C_n$ is odd if
and only if $n=0$ or $n=2m+1$ with $m=2^k-1$, that is, $n=2^{k+1}-1$.

\tools Problem~\ref{pb:kummer}; Volume~3, Exercise~1.7 (Catalan numbers); Volume~4, Chapter~2 (the equation $C=1+xC^2$).
\leads Over $\mathbb F_2$, (b) says $C(x)=\sum_kx^{2^k-1}$: a lacunary series that is a root of the quadratic $xC^2+C+1$. Such
``algebraic but sparse'' series are impossible over $\Q$: by Hadamard's gap theorem the circle of convergence of $\sum x^{2^k-1}$ is a
natural boundary, while an algebraic function has finitely many singularities. Problem~\ref{pb:thuemorse} explains what
algebraic series over $\mathbb F_p$ look like.
\end{solution}

\begin{solution}{pb:stern}
\idea The recursions for $s$ and for hyperbinary representations are the same once we look at the last digit $e_0$.

\writeup (a) Let $h(n)$ be the number of hyperbinary representations of $n$, $h(0)=1$. If $n=2m+1$ is odd then $e_0=1$, and the
remaining digits form a hyperbinary representation of $m$: $h(2m+1)=h(m)$. If $n=2m\ge2$ then $e_0\in\{0,2\}$, leaving $m$ or $m-1$:
$h(2m)=h(m)+h(m-1)$. On the other hand $u(n)=s(n+1)$ satisfies $u(2m+1)=s(2m+2)=s(m+1)=u(m)$ and
$u(2m)=s(2m+1)=s(m)+s(m+1)=u(m-1)+u(m)$, with $u(0)=1=h(0)$. By strong induction $h=u$.

(b) Expanding the product, a term $x^n$ is a choice of $e_k\in\{0,1,2\}$ in each factor with $\sum e_k2^k=n$, that is, a
hyperbinary representation of $n$. (The product converges in $\Z[[x]]$, since the $k$-th factor is $1+O(x^{2^k})$.)

(c) Let $R_k=\sum_{2^k\le n<2^{k+1}}s(n)$; $R_0=s(1)=1$. Splitting row $k+1$ into even and odd $n$,
\[
\begin{aligned}
R_{k+1}&=\sum_{2^k\le m<2^{k+1}}\bigl(s(2m)+s(2m+1)\bigr)\\
&=\sum_m\bigl(2s(m)+s(m+1)\bigr)\\
&=3R_k-s(2^k)+s(2^{k+1})=3R_k,
\end{aligned}
\]
since $s$ of a power of $2$ is $1$.

(d) $\gcd(s(0),s(1))=1$. If $\gcd(s(n),s(n+1))=1$, then $\gcd(s(2n),s(2n+1))=\gcd(s(n),s(n)+s(n+1))=1$ and
$\gcd(s(2n+1),s(2n+2))=\gcd(s(n)+s(n+1),s(n+1))=1$. By induction all consecutive pairs are coprime.

\tools Volume~1, Chapter~1 (Euclid's algorithm); Volume~4, Chapter~2 (infinite products of power series).
\leads By (d), $\frac{s(n)}{s(n+1)}$ is a fraction in lowest terms, and the sequence
\[
\frac11,\frac12,\frac21,\frac13,\frac32,\frac23,\frac31,\dots
\]
lists every positive rational exactly once: it is the Calkin--Wilf enumeration of Volume~6, Exercise~1.9 (``Recounting the
rationals''). Stern introduced the sequence in 1858 in connection with the Stern--Brocot tree of Chapter~8.
\end{solution}

\begin{solution}{pb:binarypart}
\idea Look at the smallest part: it is $1$, or every part is even and we can halve. Modulo $2$, the generating function satisfies an
equation that pins it down completely.

\writeup (a) If $n$ is odd, every partition of $n$ into powers of $2$ contains a part $1$; removing one is a bijection onto the
partitions of $n-1$ into powers of $2$, so $b(2n+1)=b(2n)$. For $2n$, the partitions with a part $1$ correspond in the same way to
those of $2n-1$, and those without have all parts even, and halving them gives the partitions of $n$. So $b(2n)=b(2n-1)+b(n)$.

(b) First proof: $b(2)=b(3)=2$. For $n\ge2$, $b(2n)=b(2n-1)+b(n)=b(2n-2)+b(n)$ is a sum of two even numbers by induction (both
$2n-2$ and $n$ are $\ge2$ and smaller than $2n$), and $b(2n+1)=b(2n)$.

Second proof: $B(x)=\sum b(n)x^n=\prod_{j\ge0}(1-x^{2^j})^{-1}$, so $(1-x)B(x)=\prod_{j\ge1}(1-x^{2^j})^{-1}=B(x^2)$. Modulo $2$, $B(x^2)=B(x)^2$, so
$(1+x)B=B^2$ in $\mathbb F_2[[x]]$, and since $B$ is invertible, $B=1+x$: all coefficients from $x^2$ on are even.

(c) Summing $b(2j)-b(2j-2)=b(j)$ for $j=1,\dots,n$ gives $b(2n)-b(0)=b(1)+\dots+b(n)$. The sequence $b$ is non-decreasing (add a part
$1$), so $b(2n)\le(n+1)b(n)$. For $n=2^k$, $b(2^k)\le\prod_{j=0}^{k-1}(2^j+1)=2^{k(k-1)/2}\prod_j(1+2^{-j})\le C\,2^{k^2/2}$. For
$2^k\le n<2^{k+1}$, monotonicity gives $\log_2b(n)\le\frac{(k+1)^2}2+O(1)\le\frac{(\log_2n)^2}2+O(\log n)$, which is the statement in
natural logarithms.

\tools Volume~4, Chapter~2 (generating functions of partitions).
\leads The upper bound has the right main term: de Bruijn (1948), following Mahler, showed that $\log b(n)\sim\frac{(\log n)^2}{2\log2}$,
with explicit lower-order terms. Churchhouse (1969) proved that $b(n)$ is never divisible by $8$ (a direct computation verifies it below $3000$), and conjectured finer congruences
for $b(2^{k+1}n)-b(2^kn)$, proved by Rødseth and Gupta.
\end{solution}

\begin{solution}{pb:arcsin}
\idea A function satisfying a linear differential equation with polynomial coefficients has Taylor coefficients satisfying a
recurrence. Here the recurrence is a ratio of quadratics, exactly the ratio of consecutive terms of $\frac{4^n}{n^2\binom{2n}n}$.

\writeup (a) $y'=\frac{2\arcsin x}{\sqrt{1-x^2}}$, so $\sqrt{1-x^2}\,y'=2\arcsin x$; differentiating, $\sqrt{1-x^2}\,y''-\frac x{\sqrt{1-x^2}}y'=\frac2{\sqrt{1-x^2}}$, and
multiply by $\sqrt{1-x^2}$.

(b) Let $S(x)=\sum_{n\ge1}c_{2n}x^{2n}$ with $c_{2n}=\frac{4^n}{2n^2\binom{2n}n}$. Since $\frac{c_{2n+2}}{c_{2n}}=\frac{(2n)^2}{(2n+1)(2n+2)}\to1$, the radius of
convergence is $1$. Writing $(1-x^2)S''-xS'=\sum_m\bigl((m+2)(m+1)c_{m+2}-m^2c_m\bigr)x^m$ (with $c_m=0$ for odd $m$ and $c_0=0$), the coefficient
of $x^0$ is $2c_2=2$ and the others vanish by the ratio just computed. So $S$ solves the same equation as $y$ with the same initial values.
For $D=y-S$, $(1-x^2)D''-xD'=0$ means $\bigl(\sqrt{1-x^2}\,D'\bigr)'=0$, so $\sqrt{1-x^2}\,D'$ is constant, equal to its value $0$ at $0$; hence $D'=0$
and $D=0$ on $]-1,1[$. At $x=\pm1$: the terms are non-negative and $c_{2n}\sim\frac{\sqrt{\pi}}2n^{-3/2}$, so the series converges at $1$,
and by monotone convergence $S(x)\to S(1)$ as $x\to1^-$, while $\arcsin^2x\to\frac{\pi^2}4$.

(c) At $x=\frac12$, $\arcsin^2\frac12=\frac{\pi^2}{36}=\frac12\sum\frac1{n^2\binom{2n}n}$. At $x=1$, $\frac{\pi^2}4=\frac12\sum\frac{4^n}{n^2\binom{2n}n}$.

\tools{Volume~4, Chapter~2 (power series) and Proposition~11.12 ($\binom{2n}n\sim\frac{4^n}{\sqrt{\pi n}}$); Volume~0, Chapter~6 (derivatives).}
\leads Euler knew (c). The same idea for $\frac1{n^4\binom{2n}n}$ gives Comtet's $\zeta(4)=\frac{36}{17}\sum_{n\ge1}\frac1{n^4\binom{2n}n}$ (1974). Integer-relation
searches by Borwein and Bradley (1997) found no formula of the same simple shape for $\zeta(5)$, which suggests that $\zeta(3)$ is special.
\end{solution}

\begin{solution}{pb:quotientlaw}
\idea The event $\lfloor1/U\rfloor=k$ is the interval $]\frac1{k+1},\frac1k]$, and inside it $\{1/U\}=\frac1U-k$. Everything is a sum over $k$ of
lengths of intervals.

\writeup (a) For $u\in[0,1]$, $P(\{1/U\}\le u)=\sum_{k\ge1}P\bigl(k\le\frac1U\le k+u\bigr)=\sum_{k\ge1}\bigl(\frac1k-\frac1{k+u}\bigr)$. The series of derivatives
$\sum\frac1{(k+u)^2}$ converges uniformly on $[0,1]$, so we may differentiate term by term (Volume~4, Chapter~2): the density is $g$. Each term
decreases in $u$; $g(0)=\zeta(2)$ and $g(1)=\zeta(2)-1$.

(b) $P(\lfloor1/U\rfloor\text{ odd})=\sum_{k\text{ odd}}\bigl(\frac1k-\frac1{k+1}\bigr)=1-\frac12+\frac13-\cdots=\log2$. And
$P(\{1/U\}<\frac12)=\sum_k\bigl(\frac1k-\frac1{k+1/2}\bigr)=2\sum_k\bigl(\frac1{2k}-\frac1{2k+1}\bigr)=2(1-\log2)$, because $\sum_k\bigl(\frac1{2k}-\frac1{2k+1}\bigr)=1-(1-\frac12+\frac13-\cdots)$.

(c) $\lfloor x/n\rfloor=k$ if and only if $\frac x{k+1}<n\le\frac xk$, which holds for $x\bigl(\frac1k-\frac1{k+1}\bigr)+O(1)$ integers $n$. Fix $K$. The $n\le x$
with $\lfloor x/n\rfloor\le K$ and odd are $x\sum_{k\le K,\,k\text{ odd}}\bigl(\frac1k-\frac1{k+1}\bigr)+O(K)$ in number, and the other $n$ with $\lfloor x/n\rfloor$ odd
satisfy $n\le\frac x{K+1}$. So the proportion lies, up to $O(\frac Kx)$, between $s_K=\sum_{k\le K\text{ odd}}(\frac1k-\frac1{k+1})$ and $s_K+\frac1{K+1}$. Let
$x\to\infty$, then $K\to\infty$: the proportion tends to $\log2$.

\tools{Volume~3, Chapter~5 (densities); Volume~4, Chapter~2 (term-by-term differentiation); Volume~0, Chapter~5 (alternating harmonic series).}
\leads The law of $\{1/U\}$ is the first step of the Gauss map $x\mapsto\{1/x\}$ of continued fractions; iterating it drives every reasonable
initial law to the Gauss measure $\frac{\mathrm dx}{(1+x)\log2}$ (Problem~\ref{pb:gausskuzmin} of Chapter~12). Note that $\log2$ appears here as the chance that
the first partial quotient of a random number is odd.
\end{solution}

\begin{solution}{pb:zigzag}
\idea The largest element splits an alternating permutation into two alternating pieces; this gives a Riccati equation, solved by $\sec+\tan$. The pole at $\frac\pi2$ then
dictates the growth, and the next pole at $-\frac{3\pi}2$ the error.

\writeup (a) Call a permutation of $m\ge2$ elements a zigzag if it is alternating starting up ($w_1<w_2>\cdots$) or starting down ($w_1>w_2<\cdots$). Complementing
($w_i\mapsto m+1-w_i$) exchanges the two kinds, so there are $2A_m$ zigzags. Let $w$ be a zigzag of $\{1,\dots,n+1\}$, $n\ge1$, and let $n+1$ be at position $k+1$. It is a peak, so
the block to its left, read from right to left, starts with a descent from $n+1$ and then goes up, down, \dots: reversed, it is an alternating permutation
(starting up) of a $k$-set, or has length at most $1$. Likewise the block to the right is alternating starting up. Conversely any choice of the $k$-set, of an alternating
arrangement of it (reversed on the left) and of an alternating arrangement of the rest gives a zigzag, of a kind determined by the parity of $k$. So
$\sum_k\binom nkA_kA_{n-k}=2A_{n+1}$.

(b) The coefficient of $\frac{z^n}{n!}$ in $f^2$ is $\sum_k\binom nkA_kA_{n-k}$, in $2f'$ it is $2A_{n+1}$; by (a) they agree for $n\ge1$, and for $n=0$, $2A_1=2=1+A_0^2$. So $2f'=1+f^2$
as formal series. The function $g=\frac{1+\sin z}{\cos z}=\sec z+\tan z$ satisfies $g(0)=1$ and
\[
2g'=2\,\frac{\cos^2z+\sin z(1+\sin z)}{\cos^2z}=\frac{2(1+\sin z)}{\cos^2z}=\frac{(1+\sin z)^2+\cos^2z}{\cos^2z}=1+g^2.
\]
The equation $2f'=1+f^2$ with $f(0)=1$ determines the Taylor coefficients one after another, so $f$ is the Taylor series of $g$, which converges on $|z|<\frac\pi2$.

(c) The zeros of $\cos$ are $\frac\pi2+j\pi$, all simple. At $z=\frac\pi2$, $1+\sin z=2$ and $(\cos)'=-1$, so $g$ has a simple pole with residue $-2$, that is principal part
$\frac{-2}{z-\pi/2}=\frac4{\pi-2z}$. At $z=-\frac\pi2$, $1+\sin z$ has a double zero, so the singularity is removable. So $g-\frac4{\pi-2z}$ is holomorphic on $|z|<\frac{3\pi}2$. By Cauchy's
estimates (Volume~4, Corollary~9.12) on a circle of radius $\rho<\frac{3\pi}2$, its coefficients are $O(\rho^{-n})$. Since $[z^n]\frac4{\pi-2z}=\frac4\pi(\frac2\pi)^n=2(\frac2\pi)^{n+1}$, the relative error is
$O\bigl((\frac\pi{2\rho})^n\bigr)$, and $\frac\pi{2\rho}$ can be any number greater than $\frac13$.

(d) $\frac{A_{10}}{10!}=\frac{50521}{3628800}=0.0139222$, and $2(\frac2\pi)^{11}=0.0139223$: the relative error is $-5.6\cdot10^{-6}$, close to $-3^{-11}$.

\tools{Volume~4, Chapter~2 (exponential generating functions), Chapter~9 (poles, Cauchy's estimates).}
\leads D\'esir\'e Andr\'e found (b) in 1879. Subtracting all the poles gives the exact formula
$\frac{A_n}{n!}=2(\frac2\pi)^{n+1}\sum_{k\ge0}(-1)^{k(n+1)}(2k+1)^{-(n+1)}$, so the error is about $(-1)^{n+1}3^{-(n+1)}$, as observed in (d). For odd $n$ the sum is
$(1-2^{-n-1})\zeta(n+1)$: the tangent numbers $A_{2m-1}$ are the coefficients behind $\zeta(2m)$, which Volume~4, Chapter~10 computes with Bernoulli numbers.
\end{solution}

\begin{solution}{pb:poissonise}
\idea Words with bounded letter counts are counted by $e_m(x)^n$, and the factor $k!$ is an integral: summing over $k$ turns the generating function into a Laplace
integral, which then has an easy asymptotic.

\writeup (a) $T>k$ means that in the first $k$ draws every value appears fewer than $m$ times. The words of length $k$ on $n$ letters with each letter used at most $m-1$
times number $k![x^k]e_m(x)^n$ (choose the multiplicities $j_1+\dots+j_n=k$, then $\frac{k!}{j_1!\cdots j_n!}$ arrangements), and each word has probability $n^{-k}$.

(b) Write $e_m(x)^n=\sum_ka_kx^k$, a polynomial. Then $\mathbb E[T]=\sum_{k\ge0}P(T>k)=\sum_ka_kk!\,n^{-k}=\sum_ka_kn^{-k}\int_0^\infty t^ke^{-t}\dd t=\int_0^\infty e^{-t}e_m(t/n)^n\dd t$.

(c) For $m=2$, $\int_0^\infty e^{-t}(1+\frac tn)^n\dd t=\sum_j\binom nj\frac{j!}{n^j}=\sum_{j=0}^n\frac{n^{\underline j}}{n^j}=1+Q(n)$. Volume~3, Exercise~2.7 is the same sum, $\sum_kq_k$.

(d) Put $t=nx$: $\mathbb E[T]=n\int_0^\infty\psi(x)^n\dd x$ with $\psi(x)=e^{-x}(1+x+\frac{x^2}2)$. Then $\psi(0)=1$ and $\psi'(x)=-\frac{x^2}2e^{-x}<0$, so $\psi$ decreases, and
$\ln\psi(x)=-\frac{x^3}6+O(x^4)$. Choose $\delta>0$ with $-\frac{x^3}6-Cx^4\le\ln\psi(x)\le-\frac{x^3}{12}$ on $[0,\delta]$. On $[\delta,\infty[$, $\psi^n\le\psi(\delta)^{n-1}\psi$ and $\int\psi=3$, which is exponentially
small. On $[0,\delta]$ put $x=(6/n)^{1/3}y$: then $n\ln\psi(x)\to-y^3$ with domination by $e^{-y^3/2}$, so
\[
n\int_0^\delta\psi^n\dd x=n\Bigl(\frac6n\Bigr)^{1/3}\int_0^\infty e^{-y^3}\dd y\,\bigl(1+o(1)\bigr)=(6n^2)^{1/3}\,\Gamma\Bigl(\frac43\Bigr)\bigl(1+o(1)\bigr),
\]
using $\int_0^\infty e^{-y^3}\dd y=\Gamma(\frac43)$ (Volume~4, Exercise~5.2).

(e) $\Gamma(\frac43)(6\cdot365^2)^{1/3}=82.9$, against the exact $88.74$. So the answer to the think box is ``about $90$''. (For pairs, $\sqrt{\pi\cdot365/2}=23.9$ against $24.62$.)

\tools{Volume~3, Exercise~2.7 (waiting for a pair); Volume~4, Chapter~2 (exponential generating functions), Chapter~5 (Gamma) and Chapter~11 (Laplace's method at an end point).}
\leads Klamkin and Newman (1967) found the general answer $\mathbb E[T]\sim(m!)^{1/m}\Gamma(1+\frac1m)\,n^{1-1/m}$ (for $m=2$ it is $\sqrt{\pi n/2}$). The relative error
decays only like $n^{-1/3}$ for $m=3$: $1\%$ at $n=10^5$. Hash tables use the same computation, with $m$ the capacity of a bucket.
\end{solution}

\begin{solution}{pb:diehit}
\idea The generating function is rational. Its only pole on the closed unit disc is at $1$, with residue $\frac1{\text{mean}}$; the other poles are farther out and only produce an
exponentially small error.

\writeup (a) Condition on the first roll $j$: the totals then hit $n$ if and only if the later totals, started afresh, hit $n-j$. So $u_n=\frac16\sum_ju_{n-j}$ for $n\ge1$, and $U=\sum u_nz^n$ satisfies
$U=1+D(z)U$.

(b) $D(1)=1$ and $D'(1)=\frac{1+\dots+6}6=\frac72\ne0$, so $1$ is a simple root. If $|z|\le1$ and $D(z)=1$, then $1=|D(z)|\le\frac16\sum|z|^j\le1$, so $|z|=1$ and all the $z^j$ have the same argument
as their sum, $1$: $z^j=1$ for $j=1,\dots,6$, so $z=1$.

(c) The polynomial $1-D$ factors as $(1-z)q(z)$ with $q(1)=D'(1)=\frac72$, and its other roots have modulus at least some $R>1$. So $V(z)=U(z)-\frac{2/7}{1-z}$ is rational without poles in
$|z|<R$, and Cauchy's estimates (Volume~4, Corollary~9.12) give $[z^n]V=O(r^{-n})$ for each $r<R$. Hence $u_n=\frac27+O(r^{-n})$. Numerically, one finds $R=1.3694$, so $\rho=0.7302$.

(d) The same argument applies to $D(z)=\sum_{j\le N}p_jz^j$: equality $D(z)=1$ on $|z|\le1$ forces $z^j=1$ for every $j$ with $p_j>0$, hence $z^g=1$ with $g$ the gcd of these $j$ (by B\'ezout, $g$ is an integer
combination of them), so $z=1$ when $g=1$; and $D'(1)=\mu$. If $g=2$, all totals are even, $u_n=0$ for odd $n$, and $z=-1$ is a second pole on the circle: $u_{2n}\to\frac2\mu$ instead.

\tools{Volume~4, Chapter~2 (generating functions) and Corollary~9.12 (Cauchy's estimates); Volume~3, Chapter~7 (first-step decompositions).}
\leads $\frac27$ is the renewal theorem: hitting points are spread at average spacing $\mu$. Erd\H os, Feller and Pollard (1949) proved $u_n\to\frac1\mu$ for any step law on the positive integers
with gcd $1$, finite support or not; without finite support there is no rational function and no exponential rate, and the proof needs a Tauberian idea.
\end{solution}

\begin{solution}{pb:flory}
\idea The first dimer cuts the row into two independent rows: a recursion with a uniform split, which becomes a first-order differential equation for the generating
function, solved by an exponential.

\writeup (a) The pair with the smallest label is filled first, and it is uniform among the $n-1$ pairs. Given that it is $\{i,i+1\}$, the labels of the other pairs are independent and
uniform on $]\ell,1]$ ($\ell$ the smallest label), and the pairs touching $i$ or $i+1$ can never be filled; the process continues independently on the rows $1,\dots,i-1$ and
$i+2,\dots,n$, each with the same rule (rescaling labels does not change their order). So $V_n=\frac1{n-1}\sum_{i=1}^{n-1}(V_{i-1}+V_{n-i-1})=\frac2{n-1}\sum_{j=0}^{n-2}V_j$.

(b) The relation $(n-1)V_n=2\sum_{j\le n-2}V_j$ also holds for $n=0,1$ (both sides vanish). Multiply by $x^n$ and sum: $xV'-V=\frac{2x^2V}{1-x}$, where $V=\sum V_nx^n$. Put $V=xW$:
$x^2W'=\frac{2x^3W}{1-x}$, so $\frac{W'}W=\frac{2x}{1-x}=-2+\frac2{1-x}$ and $W=\frac{ce^{-2x}}{(1-x)^2}$, with $c=W(0)=V_1=1$. The coefficient of $x^{n-1}$ in $\frac{e^{-2x}}{(1-x)^2}=\sum_k\frac{(-2)^kx^k}{k!}\sum_m(m+1)x^m$ is
$\sum_{k\le n-1}\frac{(-2)^k}{k!}(n-k)$.

(c) The full series is $\sum_{k\ge0}(n-k)\frac{(-2)^k}{k!}=ne^{-2}+2e^{-2}$, so $V_n-(n+2)e^{-2}=\sum_{j\ge1}j\frac{(-2)^{n+j}}{(n+j)!}$ (with $k=n+j$; up to sign). Its terms alternate in sign, and their absolute values $t_j$
satisfy $\frac{t_{j+1}}{t_j}=\frac{2(j+1)}{j(n+j+1)}\le\frac4{n+2}\le1$ for $n\ge2$, so the sum is at most $t_1=\frac{2^{n+1}}{(n+1)!}$ in absolute value. For $n=0,1$ check directly ($|0-2e^{-2}|$ and $|1-3e^{-2}|$ are less than $2$).

(d) $V=h(x)(1-x)^{-2}$ with $h(x)=xe^{-2x}$ entire, so Volume~4, Theorem~11.13 with $\alpha=2$ gives $V_n=\frac{h(1)}{\Gamma(2)}n+O(1)=e^{-2}n+O(1)$, as in (c).

\tools{Volume~4, Chapter~2 (generating functions, differential equations) and Theorem~11.13 (transfer); Volume~3, Chapter~5 (uniform labels).}
\leads Paul Flory (1939) computed this fraction for chemical groups reacting in pairs along a polymer, and Page (1959) the distribution of the vacancies. In the continuous
version, cars of length $1$ park at random on a street of length $x$, and R\'enyi (1958) showed that they fill a fraction tending to $0.7476\ldots$; a numerical calculation simulates both.
\end{solution}

\section*{Sources and further reading}

\begin{itemize}
\item H.~Wilf, \emph{generatingfunctionology}, 3rd ed., A K Peters, 2006 (free online). The source of the epigraph and of
the exponential formula in this form.
\item R.~Graham, D.~Knuth and O.~Patashnik, \emph{Concrete Mathematics}, 2nd ed., Addison-Wesley, 1994, Chapter~7.
\item P.~Flajolet and R.~Sedgewick, \emph{Analytic Combinatorics}, Cambridge University Press, 2009, Part~A (labelled and
unlabelled constructions, free online).
\item R.~Stanley, \emph{Enumerative Combinatorics}, Vol.~1, 2nd ed., Cambridge University Press, 2012, Chapter~1, and
Vol.~2, 1999, Chapter~5 (the exponential formula).
\item G.~Andrews, \emph{The Theory of Partitions}, Addison-Wesley, 1976, Chapter~1; G.~Andrews and K.~Eriksson,
\emph{Integer Partitions}, Cambridge University Press, 2004.
\item F.~Franklin, \emph{Sur le d\'eveloppement du produit infini $(1-x)(1-x^2)(1-x^3)\cdots$}, Comptes Rendus de
l'Acad\'emie des Sciences 92 (1881), 448--450.
\item G.~H.~Hardy and E.~M.~Wright, \emph{An Introduction to the Theory of Numbers}, 6th ed., Oxford University Press,
2008, Chapter~19.
\item D.~Knuth, \emph{The Art of Computer Programming}, Vol.~3, 2nd ed., Addison-Wesley, 1998, Section~5.2.2 (the analysis of
quicksort).
\item L.~Euler, \emph{Introductio in analysin infinitorum}, Lausanne, 1748, Chapter~16, where Naud\'e's question of 1740 is answered; see also G.~E.~Andrews, \emph{Euler's pioneering work on partitions}, Mathematics Magazine 56 (1983), 279--284.
\end{itemize}
